\documentclass[11pt,reqno]{amsart}

\usepackage{amsmath,amssymb,amsthm}
\usepackage{mathtools}
\usepackage[margin=1in]{geometry}
\usepackage{microtype}
\usepackage[colorlinks=true,linkcolor=blue,citecolor=blue]{hyperref}

%% ---- theorem environments ----
\newtheorem{theorem}{Theorem}[section]
\newtheorem{proposition}[theorem]{Proposition}
\newtheorem{lemma}[theorem]{Lemma}
\newtheorem{corollary}[theorem]{Corollary}
\theoremstyle{definition}
\newtheorem{definition}[theorem]{Definition}
\newtheorem{conjecture}[theorem]{Conjecture}
\newtheorem{remark}[theorem]{Remark}
\newtheorem{assumption}[theorem]{Hypothesis}

%% ---- macros ----
\newcommand{\W}{\mathcal{W}}
\newcommand{\Jc}{\mathcal{J}}
\newcommand{\Cor}{\mathcal{C}}
\newcommand{\spec}{\operatorname{spec}}
\newcommand{\specess}{\operatorname{spec_{ess}}}
\newcommand{\negind}{\operatorname{neg.ind}}
\newcommand{\ran}{\operatorname{ran}}
\newcommand{\dist}{\operatorname{dist}}
\newcommand{\tr}{\operatorname{tr}}
\newcommand{\rank}{\operatorname{rank}}
\newcommand{\R}{\mathbb{R}}
\newcommand{\C}{\mathbb{C}}
\newcommand{\N}{\mathbb{N}}
\newcommand{\Z}{\mathbb{Z}}
\newcommand{\Hil}{H}
\newcommand{\sst}{\operatorname{st}}
\newcommand{\Qaut}{Q^{\mathrm{aut}}}
\newcommand{\half}{\tfrac12}
\DeclareMathOperator{\sgn}{sgn}
\newcommand{\GAP}{\textbf{[GAP]}}

\begin{document}

\title[An Operatorial Geometry for $\zeta$]{An operatorial geometry for $\zeta$: the Weil--Toeplitz algebra and a logistic criterion for the Riemann Hypothesis}

\author{David Alejandro Trejo Pizzo}
\thanks{Independent Researcher. \texttt{dtrejopizzo@gmail.com}}

\begin{abstract}
We introduce an operatorial geometry for the Weil explicit formula. The
renormalized window forms of the explicit formula are organized into a
bounded field of finite-dimensional Hilbert spaces over the discrete base
\(\N\) (an \(\ell^\infty\)-product; no Dixmier continuity is claimed or used);
the bounded sections
form a \(C^*\)-algebra \(\W\), the sections vanishing at infinity form a closed
ideal \(\Jc\cong c_0\), and the quotient corona \(\Cor=\W/\Jc\) carries a
\emph{symbol} \(\sigma(Q)\) whose spectrum is the limit set of the window
spectra. The master quantifier (``mean over windows'' versus
``uniformly over windows'') becomes literally the distinction between the
Ces\`aro states of \(\Cor\) and its norm. Within an axiomatized spectral model
\(\mathfrak W\) we prove a \emph{trichotomy}: off-critical zeros stratify by a
resolution-dependent visibility into a fat stratum (seen by the symbol), a thin
stratum (accumulating exactly at the boundary point \(0\in\spec\sigma(Q)\)), and
a sub-resolution stratum (invisible, living in \(\Jc\)). The boundary point
\(0\) is exactly the point separating weak from strict essential positivity, and
by a sharp characterization theorem it is the unique point on which robust
finiteness of the defect depends. Independently, passing to a non-standard
normalized signature we prove an exact logistic law
\(\sigma^{(k)}=\half(1-(1-2\sigma^0)^{2^k})\) for the tensor tower of window
indices, and -- in its \emph{localized} form -- a criterion
\(\mathrm{RH}\iff\lim_k\sigma_{\mathrm{loc}}^{(k)}=0\) whose RH branch
(\(\mathrm{RH}\Rightarrow0\)) is a theorem and whose \(\neg\)RH branch
(\(\neg\mathrm{RH}\Rightarrow\half\)) is \emph{conditional} on a single, named,
purely analytic gap -- a quantitative uniformity in \(M\) of Paley--Wiener
constants in the witness construction (the Paley--Wiener uniformity problem) -- strictly narrower than
the realizability hypothesis it replaces. We are explicit about scope: this is a
reformulation, not a proof of RH and not progress toward one; the classical
wall (uniform second-order pair-correlation asymptotics) is transmuted into the
norm of the corona and into a named, non-RH, falsifiable gap; what is genuinely
new is the framework -- the Weil--Toeplitz algebra -- and the reduction of the
former realizability hypothesis under \(\neg\)RH to the strictly narrower
analytic gap above.
\end{abstract}

\maketitle

\tableofcontents
\newpage

\section{Introduction}\label{sec:intro}

\subsection{The framework in one sentence}
The Weil explicit formula defines, on a class of test functions equipped with a
notion of height, a hermitian form \(Q\) whose positivity is equivalent to the
Riemann Hypothesis (RH) \cite{Weil1952}. We package the renormalized
\emph{windows} of this form as a section of a bounded field of
finite-dimensional Hilbert spaces over \(\N\), form the associated
\(C^*\)-algebra \(\W\) of
bounded sections together with its corona \(\Cor=\W/\Jc\) (Definition
\ref{def:field}), and read the obstruction to RH inside the spectrum of the
\emph{symbol} \(\sigma(Q)\in\Cor\). This is the exact analogue, for a family,
of the passage from a Toeplitz operator to its Calkin image: the essential
spectrum is computed window-by-window in the limit.

\subsection{What is new}
Two things, and they are independent.

\emph{(A) An operatorial geometry.} The corona \(\Cor\) makes the
\emph{master quantifier} -- the gap between what arithmetic computes in mean and
what RH needs uniformly -- into a structural feature: membership in \(c_0\)
versus \(\ell^\infty\), equivalently the difference between the Ces\`aro states
of \(\Cor\) and its norm (Section \ref{sec:trichotomy}, Section
\ref{sec:logistic}). Inside this geometry we prove a trichotomy theorem
(Theorems \ref{thm:fat}--\ref{thm:nogo}): off-critical zeros do not split
dichotomously into ``visible'' and ``invisible''; they stratify by visibility,
and the intermediate stratum lives exactly at the single boundary point
\(0\in\spec\sigma(Q)\). That point is, by a sharp characterization (Theorem
\ref{thm:exact}), the unique point on which the Fredholm half of the architecture
depends.

\emph{(B) A criterion whose RH branch is a theorem.} In normalized
non-standard coordinates the tensor tower of window indices obeys an exact
logistic recursion (Theorem \ref{thm:logistic}), and in its localized version we
obtain
\[
\mathrm{RH}\;\Longleftrightarrow\;\lim_{k\to\infty}\sigma_{\mathrm{loc}}^{(k)}=0
\]
with the RH branch a theorem (compressed Weil positivity, Proposition
\ref{prop:RHbranch}) and the \(\neg\)RH branch \emph{conditional} on a single,
named, purely analytic gap: a quantitative uniformity in \(M\) of the
Paley--Wiener constants of the witness construction (Hypothesis \ref{hyp:pw},
the Paley--Wiener uniformity problem; Theorem \ref{thm:realiz}). This gap is strictly narrower than the
realizability hypothesis it replaces (Remark \ref{rem:key}), but it is open: the
\(\neg\)RH branch is \emph{not} claimed as a theorem here.

\subsection{Honest scope, stated up front}\label{subsec:honest-up-front}
We are categorical about this and devote Section \ref{sec:honest} to it. (a) We
do \emph{not} prove RH and we claim no progress toward a proof: the equivalence
of (B) is a reformulation. (b) The ``computability of the symbol'' -- whether
\(\sigma(Q)\) equals the autonomous arithmetic prediction -- is the classical
wall (uniform second-order pair-correlation asymptotics) transmuted into the
norm topology of the corona; it is not resolved (Section \ref{sec:honest},
Theorem \ref{thm:symbolcompute}). (c) What is genuinely new is the framework and
the reduction of the former realizability hypothesis under \(\neg\)RH to the
strictly narrower analytic gap of the Paley--Wiener uniformity problem (Hypothesis \ref{hyp:pw}); that
gap is open, and the \(\neg\)RH branch is conditional on it. Every result
below is labelled: \emph{Theorem} (complete proof here), \emph{Proposition} /
\emph{Lemma} (auxiliary, complete proof here), or \emph{Conjecture} / \emph{Open}
/ \(\GAP\) (declared, not proved). When a step rests on a structural input from
the author's earlier work we name it as a cited hypothesis, never assume it
tacitly.

\subsection{Relation to the author's earlier work}
This paper is self-contained: all proofs are given here. It builds on three
earlier works of the author, cited as previous work and never re-proved: the
\(\kappa=2m\) bridge identification of the Weil form on a Pontryagin space
\cite{P35}; the trace criterion and the autonomy of the value of \(Q\) computed
from primes without zero positions \cite{P39}; and the spectral structure of the
windowed Weil form (online sea \(\approx\) identity, each visible off-critical
quadruple a hyperbolic plane) \cite{P40}. The precise statements we import are
collected as Hypotheses \ref{hyp:blocks}--\ref{hyp:spectral} and tabulated in
Section \ref{sec:honest}.

\subsection*{Notation}
\(H\) denotes a complex separable Hilbert space; \(\mathcal B(H)\) the bounded
operators; \(G=G^*\) a bounded self-adjoint operator with spectral measure
\(E_G\). A hermitian form is sesquilinear, linear in the first variable, with
\(Q(y,x)=\overline{Q(x,y)}\). For a subspace \(W\), \(\negind(Q|_W)\) is the
supremum of \(\dim N\) over subspaces \(N\subseteq W\) with \(Q(x,x)<0\) for all
\(x\in N\setminus\{0\}\), valued in \(\N\cup\{\infty\}\). A nontrivial zero is
written \(\rho=\beta+i\gamma\), \(\beta\in(0,1)\); a zero is \emph{off-critical}
if \(\beta\neq\half\); a quadruple is \(\{\half\pm\delta\pm i\gamma\}\).

\section{Defect spaces and the defect \texorpdfstring{$\delta$}{delta}}
\label{sec:defect}

We isolate the minimal language of index needed below. The defect of a hermitian
form is the dimension of a maximal negative subspace; it is the
polarization-theoretic analogue of the Fredholm index.

\begin{definition}[defect]\label{def:defect}
For a hermitian form \(Q\) on a complex vector space \(H\) with a marked
subspace \(P\), the \emph{defect} is
\(\delta(Q,P):=\negind(Q|_P)\in\N\cup\{\infty\}\). When \(H\) is Hilbert,
\(P=H\), and \(Q(x,y)=\langle Gx,y\rangle\) with \(G=G^*\) bounded, we write
\(\delta(G)\). The form is \emph{polarized} if \(\delta=0\).
\end{definition}

\begin{lemma}[the defect is spectral]\label{lem:defect-spectral}
\(\delta(G)=\dim\ran E_G(-\infty,0)\).
\end{lemma}

\begin{proof}
On \(N_-:=\ran E_G(-\infty,0)\) one has
\(\langle Gx,x\rangle=\int_{(-\infty,0)}\lambda\,d\|E_\lambda x\|^2<0\) for
\(x\neq0\); thus \(N_-\) is negative and \(\delta\geq\dim N_-\). Conversely, if
\(N\) is negative then \(E_G(-\infty,0)|_N\) is injective: if
\(E_G(-\infty,0)x=0\) then \(x\in\ran E_G[0,\infty)\), whence
\(\langle Gx,x\rangle\geq0\), impossible for \(x\in N\setminus\{0\}\). So
\(\dim N\leq\dim N_-\).
\end{proof}

\subsection{Arithmetic of the defect}
The defect is monotone, additive, and tensor-multiplicative in a precise sense.
These are load-bearing only through the \emph{tensor identity}, which drives the
logistic law of Section \ref{sec:logistic}; we re-prove what we use.

\begin{lemma}[monotonicity and additivity]\label{lem:mono-add}
\textup{(i)} If \(f\colon(H_1,Q_1,P_1)\to(H_2,Q_2,P_2)\) is linear, isometric
\((Q_2(fx,fy)=Q_1(x,y))\) with \(f(P_1)\subseteq P_2\), then
\(\delta(Q_1,P_1)\leq\delta(Q_2,P_2)\). \textup{(ii)}
\(\delta(Q_1\oplus Q_2)=\delta(Q_1)+\delta(Q_2)\), with no dimension or
nondegeneracy hypothesis.
\end{lemma}

\begin{proof}
(i) For \(N\subseteq P_1\) negative, \(f|_N\) is injective (\(fx=0\Rightarrow
Q_1(x,x)=Q_2(fx,fx)=0\Rightarrow x=0\)) and \(f(N)\subseteq P_2\) is negative by
isometry. (ii) ``\(\geq\)'' is direct sum of witnesses. ``\(\leq\)'': if some
\(\delta(Q_i)=\infty\) there is nothing to prove; if both are finite, choose
maximal negatives \(N_i\) with \(P_i=N_i\oplus R_i\) (\(Q_i\geq0\) on \(R_i\):
if \(w\in R_i\) had \(Q_i(w,w)<0\) then \(N_i\oplus\C w\) would be negative of
larger dimension). For \(N\subseteq P_1\oplus P_2\) negative, the projection
\(\pi\colon N\to N_1\oplus N_2\) along \(R_1\oplus R_2\) has trivial kernel
(\(\ker\pi=N\cap(R_1\oplus R_2)\) is both \(\geq0\) and \(<0\)), so
\(\dim N\leq\delta(Q_1)+\delta(Q_2)\).
\end{proof}

\begin{lemma}[Whitney formula for the defect]\label{lem:whitney}
For finite-dimensional nondegenerate hermitian forms of signatures
\((p_1,q_1)\), \((p_2,q_2)\) (with \(q_i=\delta\)),
\[
\sgn(Q_1\otimes Q_2)=(p_1p_2+q_1q_2,\;p_1q_2+q_1p_2),
\qquad
\delta(Q_1\otimes Q_2)=p_1q_2+q_1p_2 .
\]
\end{lemma}

\begin{proof}
Take \(Q_i\)-orthonormal bases with signs \(\varepsilon^{(i)}_a\in\{\pm1\}\)
(Sylvester's law of inertia \cite{Gantmacher}). The tensor basis is \(Q_1\otimes Q_2\)-orthogonal with signs
\(\varepsilon^{(1)}_a\varepsilon^{(2)}_b\); the sign is \(-1\) exactly when the
two factors disagree, giving \(p_1q_2+q_1p_2\) negative directions and a
diagonal orthogonal basis with no isotropic vectors (nondegeneracy).
\end{proof}

\begin{corollary}[reduced signature is the only multiplicative character]
\label{cor:character}
With \(d=p+q\), \(s=p-q\): \(d(Q_1\otimes Q_2)=d_1d_2\) and
\(s(Q_1\otimes Q_2)=s_1s_2\), and \(\delta=(d-s)/2\). In particular
\(\delta\) is not multiplicative; the Grothendieck ring of finite nondegenerate
forms under \(\oplus,\otimes\) is \(\Z[\varepsilon]/(\varepsilon^2-1)\) via
\(p+q\varepsilon\), whose only ring homomorphisms to an integral domain are
\(\varepsilon\mapsto\pm1\), i.e.\ \(d\) and \(s\).
\end{corollary}

\begin{proof}
\(s_1s_2=(p_1-q_1)(p_2-q_2)=(p_1p_2+q_1q_2)-(p_1q_2+q_1p_2)\) by Lemma
\ref{lem:whitney}; \(d_1d_2=(p_1+q_1)(p_2+q_2)\) is total dimension. A ring
homomorphism to an integral domain sends \(\varepsilon\) to a root of
\(\varepsilon^2=1\).
\end{proof}

\subsection{The exact characterization of robust finiteness}
This is the first load-bearing theorem. Essential positivity comes in two
strengths, and only the strict one is stable under the compact ideal; the
distinction is the whole content.

\begin{definition}[two essential positivities]\label{def:two-pos}
\(G=G^*\) is \emph{weakly essentially positive} if
\(\specess(G)\subseteq[0,\infty)\), and \emph{strictly essentially positive} if
\(\specess(G)\subseteq[c,\infty)\) for some \(c>0\). Since \(\specess(G)\) is
compact, the latter is equivalent to \(\specess(G)\subseteq(0,\infty)\).
\end{definition}

\begin{theorem}[exact characterization of robust finiteness of the defect]
\label{thm:exact}
For \(G=G^*\in\mathcal B(H)\) the following are equivalent:
\begin{enumerate}
\item[\textup{(i)}] \(G\) is strictly essentially positive;
\item[\textup{(ii)}] \(\delta(G+K)<\infty\) for \emph{every} compact
self-adjoint \(K\);
\item[\textup{(iii)}] there is \(\varepsilon>0\) with
\(\delta(G-\varepsilon I)<\infty\).
\end{enumerate}
\end{theorem}

\begin{proof}
\emph{(i)\(\Rightarrow\)(ii).} By Weyl's theorem on stability of the essential
spectrum under compact perturbations \cite[Thm.\ IV.5.35]{Kato},
\(\specess(G+K)=\specess(G)\subseteq[c,\infty)\). Hence
\(\spec(G+K)\cap(-\infty,0]\) is bounded and discrete (its accumulation points
would lie in \(\specess(G+K)\), disjoint from \((-\infty,0]\)), so finite, of
eigenvalues of finite multiplicity, and \(\delta(G+K)<\infty\) by Lemma
\ref{lem:defect-spectral}.

\emph{(iii)\(\Leftrightarrow\)(i).} \(\delta(G-\varepsilon I)=\dim\ran
E_G(-\infty,\varepsilon)\). If this is finite, \(\spec(G)\cap(-\infty,\varepsilon)\)
is finitely many eigenvalues of finite multiplicity, so
\(\specess(G)\subseteq[\varepsilon,\infty)\): strict. Conversely, if
\(\specess(G)\subseteq[c,\infty)\), take \(\varepsilon=c/2\); the same argument
gives \(\dim\ran E_G(-\infty,c/2)<\infty\).

\emph{(ii)\(\Rightarrow\)(i), by contraposition.} Suppose strictness fails:
there is \(\mu\in\specess(G)\) with \(\mu\leq0\). We construct a compact
self-adjoint \(K\) with \(\delta(G+K)=\infty\).

\emph{Step 1 (orthonormal Weyl sequence with summable error).} For each \(n\),
\(\dim\ran E_G(\mu-8^{-n},\mu+8^{-n})=\infty\) (this is the defining property of
the essential spectrum for self-adjoint operators: if it were finite for some
\(n\), \(\mu\) would be an isolated eigenvalue of finite multiplicity, hence not
in \(\specess\) \cite[\S XIII.4]{ReedSimon}). Choose inductively
\(e_n\in\ran E_G(\mu-8^{-n},\mu+8^{-n})\) of unit norm orthogonal to
\(e_1,\dots,e_{n-1}\) (possible: the range is infinite-dimensional and finitely
many orthogonality constraints). Then \((e_n)\) is orthonormal and
\[
\|(G-\mu)e_n\|^2=\int_{(\mu-8^{-n},\mu+8^{-n})}|\lambda-\mu|^2\,d\|E_\lambda e_n\|^2
\leq 64^{-n},\qquad\text{i.e.}\quad\|(G-\mu)e_n\|\leq 8^{-n}.
\]

\emph{Step 2 (Gram matrix of \(G\)).} For \(n\neq m\), by self-adjointness,
\(|\langle Ge_n,e_m\rangle|=|\langle(G-\mu)e_n,e_m\rangle|\leq 8^{-n}\) and
\(=|\langle e_n,(G-\mu)e_m\rangle|\leq 8^{-m}\), so
\(|\langle Ge_n,e_m\rangle|\leq 8^{-\max(n,m)}\); and
\(\langle Ge_n,e_n\rangle\leq\mu+8^{-n}\leq8^{-n}\) (as \(\mu\leq0\)).

\emph{Step 3 (the perturbation).} Set
\(K:=-\sum_{n\geq1}2^{-n}\langle\cdot,e_n\rangle e_n\): compact (norm limit of
finite-rank), self-adjoint, \(\|K\|=\half\). For
\(x=\sum_{n\geq n_0}\xi_n e_n\) finitely supported in \(\{n\geq n_0\}\),
\[
\langle(G+K)x,x\rangle\leq\sum_n(8^{-n}-2^{-n})|\xi_n|^2
+\sum_{n\neq m}8^{-\max(n,m)}|\xi_n||\xi_m|.
\]
The cross term is bounded by
\(\sum_{n\neq m}8^{-\max(n,m)}\tfrac{|\xi_n|^2+|\xi_m|^2}{2}
\leq\sum_n|\xi_n|^2 r_n\) with
\(r_n=\sum_{m\neq n}8^{-\max(n,m)}\leq n8^{-n}+\sum_{m>n}8^{-m}\leq(n+1)8^{-n}\).
Hence
\(\langle(G+K)x,x\rangle\leq\sum_n\bigl(8^{-n}+(n+1)8^{-n}-2^{-n}\bigr)|\xi_n|^2\).
Since \((n+2)8^{-n}<2^{-n}\) for all \(n\geq4\), taking \(n_0=4\) makes every
coefficient negative: \(G+K\) is negative definite on the
infinite-dimensional \(V_0:=\overline{\operatorname{span}}\{e_n:n\geq4\}\), so
\(\delta(G+K)=\infty\) by Lemma \ref{lem:defect-spectral}.
\end{proof}

\begin{remark}[the design lesson]\label{rem:design}
The only reading of ``essential positivity \(\Rightarrow\) finite defect'' that
survives the compact ideal -- and the essential spectrum lives modulo that ideal
by definition -- is the \emph{strict} one. Any certificate must produce a
\emph{gap} \(c>0\), not merely a sign. The operator \(G=\operatorname{diag}(-1/n)\)
has \(\specess(G)=\{0\}\subseteq[0,\infty)\) (weakly positive) yet
\(\delta=\infty\): by Theorem \ref{thm:exact}, this is the generic shape of
gaplessness. The (i)\(\Rightarrow\)(ii) direction is the form in which earlier
work \cite{P40} used Weyl's theorem; the contrapositive (ii)\(\Rightarrow\)(i),
proved above, is what makes the characterization \emph{exact} and is new here.
\end{remark}

\begin{lemma}[gap detector]\label{lem:gap-detector}
If there is an orthonormal sequence \((u_n)\) with
\(\liminf_n\langle Gu_n,u_n\rangle\leq\ell\), then
\(\inf\specess(G)\leq\ell\).
\end{lemma}

\begin{proof}
Suppose \(\specess(G)\subseteq[c,\infty)\) with \(c>\ell\); fix
\(\varepsilon\in(0,c-\ell)\) and \(E:=E_G(-\infty,c-\varepsilon)\), finite-rank
by Theorem \ref{thm:exact}(iii)\(\Leftrightarrow\)(i). With
\(m_0:=\inf\spec(G)>-\infty\),
\[
\langle Gu,u\rangle\geq m_0\|Eu\|^2+(c-\varepsilon)\|(1-E)u\|^2
=(c-\varepsilon)\|u\|^2-(c-\varepsilon-m_0)\|Eu\|^2 .
\]
Since \(u_n\rightharpoonup0\) and \(E\) is finite-rank, \(\|Eu_n\|\to0\), so
\(\liminf\langle Gu_n,u_n\rangle\geq c-\varepsilon>\ell\), a contradiction.
\end{proof}

\section{The Weil--Toeplitz algebra}\label{sec:algebra}

\subsection{The field of windows and the corona}

\begin{definition}[field of windows; \(\W\), \(\Jc\), the symbol]\label{def:field}
A \emph{field of windows} is a sequence
\(\{(H_n,Q_n)\}_{n\geq2}\) with \(H_n\) a finite-dimensional Hilbert space
\((\dim H_n=d_n<\infty)\), \(Q_n=Q_n^*\in\mathcal B(H_n)\), and
\(\sup_n\|Q_n\|<\infty\). Set
\[
\W:=\Bigl\{A=(A_n)_n:A_n\in\mathcal B(H_n),\ \|A\|:=\sup_n\|A_n\|<\infty\Bigr\},
\qquad
\Jc:=\{A\in\W:\|A_n\|\to0\}.
\]
With pointwise operations, \(\W\) is a \(C^*\)-algebra (the \(\ell^\infty\)
product of the field) and \(\Jc\) is a closed two-sided ideal (the
\emph{compact ideal of the field}: it is the closure of the finitely supported
sections). The \emph{corona} is \(\Cor:=\W/\Jc\) and the \emph{symbol} of a
section is its class \(\sigma(A):=A+\Jc\in\Cor\).
\end{definition}

\begin{remark}
In the block-diagonal operator \(\bigoplus_n A_n\) on \(\bigoplus_n H_n\), a
section lies in \(\Jc\) iff the operator is compact; \(\Cor\) is the image of
the section algebra in the Calkin algebra of \(\bigoplus_n H_n\). The corona is
literally ``the family modulo compacts''. The name Weil--Toeplitz reflects the
analogy with the Toeplitz symbol: the essential spectrum of \(T_\phi\) is the
range of the symbol -- the ``limits of the windows'' of the operator
\cite{BottcherSilbermann,RabinovichRochSilbermann}.
\end{remark}

\begin{lemma}[norm of the symbol]\label{lem:symbol-norm}
\(\|\sigma(A)\|_{\Cor}=\limsup_n\|A_n\|\).
\end{lemma}

\begin{proof}
(\(\leq\)) For fixed \(N\) the truncation
\(A^{(N)}:=(A_2,\dots,A_N,0,0,\dots)\in\Jc\), so
\(\|\sigma(A)\|\leq\|A-A^{(N)}\|=\sup_{n>N}\|A_n\|\to\limsup_n\|A_n\|\). (\(\geq\))
For every \(J\in\Jc\),
\(\|A+J\|=\sup_n\|A_n+J_n\|\geq\limsup_n\|A_n+J_n\|=\limsup_n\|A_n\|\) since
\(\|J_n\|\to0\); take the infimum over \(J\).
\end{proof}

\begin{lemma}[symbol spectrum is the limit set of window spectra]
\label{lem:symbol-spectrum}
For \(A=(A_n)\in\W\) with each \(A_n\) self-adjoint,
\[
\spec_{\Cor}(\sigma(A))=\Lambda(A):=\bigl\{\lambda\in\R:\exists\,n_k\to\infty,\
\lambda_k\in\spec(A_{n_k}),\ \lambda_k\to\lambda\bigr\}.
\]
\end{lemma}

\begin{proof}
\(\sigma(A)\) is self-adjoint, so its spectrum is real; \(\Lambda(A)\) is closed
and bounded.

(\(\subseteq\) via contrapositive) If \(\sigma(A)-\lambda\) is invertible in
\(\Cor\), there is \(B\in\W\) with \((A-\lambda)B=1+J\), \(B(A-\lambda)=1+J'\),
\(J,J'\in\Jc\). For large \(n\), \(\|J_n\|,\|J'_n\|\leq\half\), so
\((A_n-\lambda)B_n\) and \(B_n(A_n-\lambda)\) are invertible, whence
\(A_n-\lambda\) is invertible with \(\|(A_n-\lambda)^{-1}\|\leq2\|B\|\). For
self-adjoint \(A_n\),
\(\|(A_n-\lambda)^{-1}\|=\dist(\lambda,\spec A_n)^{-1}\), so
\(\dist(\lambda,\spec A_n)\geq(2\|B\|)^{-1}\) for all large \(n\): no
\(\lambda_k\in\spec(A_{n_k})\) can converge to \(\lambda\), i.e.\
\(\lambda\notin\Lambda(A)\).

(\(\supseteq\) via contrapositive) If \(\lambda\notin\Lambda(A)\), there are
\(\varepsilon>0\) and \(N\) with \(\dist(\lambda,\spec A_n)\geq\varepsilon\) for
all \(n\geq N\) (else a violating subsequence would, by boundedness of the
spectra, produce a point of \(\Lambda\) within \(\varepsilon\) of \(\lambda\)).
Put \(B_n:=(A_n-\lambda)^{-1}\) for \(n\geq N\) (\(\|B_n\|\leq\varepsilon^{-1}\))
and \(B_n:=0\) for \(n<N\); then \(B\in\W\) and \((A-\lambda)B-1\),
\(B(A-\lambda)-1\) are finitely supported, hence in \(\Jc\), so \(\sigma(B)\)
inverts \(\sigma(A)-\lambda\).
\end{proof}

\begin{corollary}[positivity of the symbol is the uniform quantifier]
\label{cor:master-quantifier}
For \(A=A^*\):
\begin{itemize}
\item \(\spec(\sigma(A))\subseteq[0,\infty)\) \emph{(weak symbol positivity)}
\(\iff\) \(\min\spec(A_n)\geq-\varepsilon_n\) with \(\varepsilon_n\to0\);
\item \(\spec(\sigma(A))\subseteq[c,\infty)\), \(c>0\) \emph{(strict symbol
positivity)} \(\iff\) for all \(\varepsilon>0\) there is \(N\) with
\(\spec(A_n)\subseteq(c-\varepsilon,\infty)\) for all \(n\geq N\) -- a gap
\emph{uniform over windows, eventually}.
\end{itemize}
\end{corollary}

\begin{proof}
Read Lemma \ref{lem:symbol-spectrum} twice.
\end{proof}

This is the structural observation that orients the paper: the master
quantifier (mean over windows versus uniformly over windows) becomes membership
in \(c_0\) versus \(\ell^\infty\). Corollary \ref{cor:master-quantifier} is the
statement ``uniform in every window'' written as positivity of an element of a
\(C^*\)-algebra.

\subsection{The field of Weil}
We now describe the field carrying the Weil form. The structural inputs we
import from earlier work are isolated as hypotheses.

\begin{definition}[the Weil field at a resolution calendar]\label{def:weil-field}
Fix \(\Delta\in(0,1]\) and a nondecreasing \emph{resolution calendar}
\(a\colon[2,\infty)\to[1,\infty)\) (the log-support budget of the tests per
window; the natural choice is \(a(T)\asymp\log T\), the resolution scale of
individual zeros). For each \(n\): \(H_n\) is a finite-dimensional subspace of
the test class, concentrated in the window \(|\operatorname{Im}s-n|\leq\Delta\)
with multiplicative support \(\subseteq[e^{-a(n)},e^{a(n)}]\), endowed with the
\emph{online-sea normalization} (the inner product is the spectral Gram form of
the online zeros of the window, normalized to mass \(1\) per zero); and \(Q_n\)
is the compression of the Weil form to \(H_n\) in that normalization. The
\emph{autonomous} section \(\Qaut=(\Qaut_n)\) is the polar/archimedean-plus-near-
prime prediction of each window, computable without zero positions
\cite{P39,P40}.
\end{definition}

The following are imported, not proved here.

\begin{assumption}[block structure {\cite{P40}}]\label{hyp:blocks}
Each \(Q_n\) decomposes, up to controlled couplings, as an online-sea block
(\(\approx\) identity after normalization) plus one hyperbolic plane
\(\left(\begin{smallmatrix}0&\theta\\\theta&0\end{smallmatrix}\right)\) for each
visible off-critical quadruple in the window.
\end{assumption}

\begin{assumption}[autonomy of the value {\cite{P39,P35}}]\label{hyp:autonomy}
For tests \(f,g\) of the class, \(Q_n(f,g)\) is computable from the arithmetic
blocks of the explicit formula (primes, pole, gamma factor) without any zero
position.
\end{assumption}

\begin{assumption}[spectral realization {\cite{P35}}]\label{hyp:spectral}
The Weil form is realized on a Pontryagin space \((\mathcal K,Q)\) with
\(\kappa(Q)=2m\), where \(m\) is the number of off-critical quadruples;
\(\mathrm{RH}\iff m=0\). \emph{(The identification of \(\kappa\) with the
negative index of the scaling generator is a conjecture of \cite{P35}; we do not
use it.)}
\end{assumption}

\begin{remark}[honest status of the field]\label{rem:field-status}
That the compressions exist with the stated normalization, and that the section
is bounded, uses the realizability apparatus of \cite{P40}. The uniform
boundedness \(\sup_n\|Q_n\|<\infty\) in the online-sea normalization requires
control of the per-window pair-correlation corrections; we treat this as an
axiom \textbf{(S1)} of the model \(\mathfrak W\) below
(\(\GAP\): per-window pair-correlation boundedness, not proved here), inside
which all theorems of Section \ref{sec:trichotomy} are proved with complete
rigor. This separates, as the contract demands, \emph{theorems} (in the
axiomatized model) from the \emph{bridge} (the genuine Weil instance, Section
\ref{sec:honest}).
\end{remark}

\section{The spectral model and the trichotomy}\label{sec:trichotomy}

\subsection{The model}
We axiomatize the spectral structure of the windowed Weil form, distilled from
Hypotheses \ref{hyp:blocks}--\ref{hyp:spectral}.

\begin{definition}[configuration and visibility]\label{def:config}
A \emph{configuration} is an at most countable set
\(Z=\{(\delta_j,\gamma_j)\}_{j\in J}\) with \(\delta_j\in(0,\half)\),
\(\gamma_j\to\infty\) if \(|J|=\infty\), and finitely many \(j\) per window;
\(m:=|J|\). (\(Z=\emptyset\) is ``RH''.) The \emph{visibility} of quadruple
\(j\) at calendar \(a\) is \(\theta_j:=\theta(\delta_j,\gamma_j)\in[0,1]\),
where \(\theta\) is a fixed function with
\begin{itemize}
\item[\textup{(V1)}] \(\theta\) continuous, nondecreasing in \(\delta\), and
\(\theta(\delta,\gamma)\to0\) as \(\delta\,a(\gamma)\to0\);
\item[\textup{(V2)}] there is \(c_0>0\) with \(\theta(\delta,\gamma)\geq c_0\)
when \(\delta\,a(\gamma)\geq1\).
\end{itemize}
Quadruples with \(\delta_j\,a(\gamma_j)\geq1\) are \emph{supra-resolution (fat)};
with \(\delta_j\,a(\gamma_j)<1\), \emph{sub-resolution (thin)}.
\end{definition}

\begin{remark}[origin of \(\theta\)]\label{rem:theta}
(V1)--(V2) are the correct properties, justified by a Paley--Wiener estimate: a
test with multiplicative support in \([e^{-a},e^a]\) has
\(\hat f(s)=\int_{-a}^{a}F(u)e^{(s-1/2)u}\,du\), so over a band of width
\(2\delta\),
\[
\hat f(\tfrac12+\delta+i\gamma)\,\overline{\hat f(\tfrac12-\delta+i\gamma)}
=\bigl|\hat f(\tfrac12+i\gamma)\bigr|^2
+O\bigl(\delta\,a\,e^{a/2}\|F\|_{L^1}\sup_{\text{band}}|\hat f|\bigr).
\]
Thus a quadruple with \(\delta a\to0\) differs from an online pair by
\(O(\delta a)\to0\) per unit mass (V1); for \(\delta a\geq1\) an aligned test
realizes a sign opposition of size \(\geq c_0\) (V2), using closure of the class
under multiplication by Mellin polynomials \cite{P40}. The exact exponent in the
transition \(\delta a\asymp1\) is not load-bearing (\(\GAP\), declared,
non-load-bearing): no theorem below uses more than (V1)--(V2).
\end{remark}

\begin{definition}[the model \(\mathfrak W(Z)\)]\label{def:model}
Given a configuration \(Z\) and calendar \(a\), set
\[
H_n:=\C^{d_n}\oplus\bigoplus_{j\in J_n^{\mathrm{vis}}}\C^2,\qquad
Q_n:=(I_{d_n}+R_n)\oplus\bigoplus_{j\in J_n^{\mathrm{vis}}}
\begin{pmatrix}0&\theta_j\\\theta_j&0\end{pmatrix}+C_n,
\]
where \(d_n\asymp\Delta\log n\),
\(J_n^{\mathrm{vis}}:=\{j:|\gamma_j-n|\leq\Delta,\ \theta_j>0\}\), and
\begin{itemize}
\item[\textup{(S1)}] \(\|R_n\|\leq r<\half\) for all \(n\) \emph{(sea
corrections; \(\GAP\) for the Weil instance, Remark \ref{rem:field-status})};
\item[\textup{(S2)}] the hyperbolic blocks carry weight \(\theta_j\) with
(V1)--(V2) \emph{(Remark \ref{rem:theta} for the Weil instance)};
\item[\textup{(S3)}] \(\|C_n\|\leq\omega_n\) with \(\omega_n\to0\)
\emph{(cross-couplings between separated windows; \(\GAP\) for the Weil
instance, inherited from correlation decay)}.
\end{itemize}
The autonomous section is \(\Qaut_n:=(I_{d_n}+R_n)\oplus
I_{2|J_n^{\mathrm{vis}}|}\). Sub-resolution quadruples with \(\theta_j=0\) do not
appear in \(H_n\); their entire effect is absorbed into \(R_n/C_n\).
\end{definition}

All theorems in this section are \emph{unconditional inside} \(\mathfrak W\);
their transport to \(\zeta\) carries the statuses of (S1)--(S3) and is discussed
in Section \ref{sec:honest}.

\subsection{The stratification theorems}
Throughout, \(\sigma(Q)\in\Cor\), \(\Lambda(Q)=\spec(\sigma(Q))\), and
\(\vartheta_n:=\max\{\theta_j:j\in J_n^{\mathrm{vis}}\}\) (\(:=0\) if empty).

\begin{lemma}[per-window spectrum]\label{lem:window-spec}
For all \(n\),
\(\spec(Q_n)\subseteq\bigl([1-r,1+r]\cup\{\pm\theta_j\}_{j\in J_n^{\mathrm{vis}}}
\bigr)+[-\omega_n,\omega_n]\). In particular
\(\lambda_{\min}(Q_n)\geq-\vartheta_n-\omega_n\); and if
\(J_n^{\mathrm{vis}}\neq\emptyset\),
\(\lambda_{\min}(Q_n)\leq-\vartheta_n+\omega_n\).
\end{lemma}

\begin{proof}
Without \(C_n\), \(Q_n\) is a direct sum: the sea block has spectrum in
\([1-r,1+r]\), and each
\(\left(\begin{smallmatrix}0&\theta\\\theta&0\end{smallmatrix}\right)\) has
eigenvalues \(\pm\theta\). The self-adjoint perturbation \(C_n\) moves each
eigenvalue by at most \(\|C_n\|\leq\omega_n\) (finite-dimensional Weyl
inequality \cite{Bhatia}). The upper bound: the vector \(v=(1,-1)/\sqrt2\) of
the \(\vartheta_n\)-block gives
\(\lambda_{\min}(Q_n)\leq\langle Q_nv,v\rangle\leq-\vartheta_n+\omega_n\).
\end{proof}

\begin{theorem}[fat \(\Rightarrow\) the symbol sees them]\label{thm:fat}
If there are \(c>0\) and a subsequence \(j_k\) with \(\theta_{j_k}\geq c\) and
\(\gamma_{j_k}\to\infty\), then
\(\spec(\sigma(Q))\cap(-\infty,-c/2]\neq\emptyset\). In particular the symbol is
not even weakly positive: the negativity of supra-resolution quadruples
\emph{survives the quotient by the compact ideal}.
\end{theorem}

\begin{proof}
Let \(n_k\) be the window of \(\gamma_{j_k}\) (\(n_k\to\infty\)). By Lemma
\ref{lem:window-spec}, \(\lambda_{\min}(Q_{n_k})\leq-c+\omega_{n_k}\). Since
\(\omega_n\to0\) and the spectra are uniformly bounded, a subsequence of
\(\lambda_{\min}(Q_{n_k})\) converges to a point \(\leq-c\), which lies in
\(\Lambda(Q)\) by Lemma \ref{lem:symbol-spectrum}.
\end{proof}

\begin{theorem}[thin visible \(\Rightarrow\) negatives live exactly at the
boundary \(0\)]\label{thm:thin}
Suppose \(m=\infty\), \(J_n^{\mathrm{vis}}\neq\emptyset\) for infinitely many
\(n\), and \(\vartheta_n\to0\). Then:
\begin{enumerate}
\item[\textup{(i)}] \(0\in\spec(\sigma(Q))\); if moreover \(\vartheta_n>\omega_n\)
for infinitely many populated windows, it is attained as a limit of genuinely
\emph{negative} \(\lambda_{\min}(Q_n)\) over those windows;
\item[\textup{(ii)}] \(\spec(\sigma(Q))\subseteq\{0\}\cup[1-r,1+r]\): weak symbol
positivity \emph{holds};
\item[\textup{(iii)}] strict symbol positivity \emph{fails}.
\end{enumerate}
The negatives accumulate exactly at the single point \(0\) -- the point
separating weak from strict positivity, the very point of Theorem
\ref{thm:exact}.
\end{theorem}

\begin{proof}
(i) Over the infinitely many populated windows,
\(\lambda_{\min}(Q_n)\in[-\vartheta_n-\omega_n,-\vartheta_n+\omega_n]\to0\)
(Lemma \ref{lem:window-spec}, \(\vartheta_n,\omega_n\to0\)); when
\(\vartheta_n>\omega_n\) these are genuinely negative, and the limit \(0\) lies
in \(\Lambda(Q)\) (Lemma \ref{lem:symbol-spectrum}). (ii) Every eigenvalue of
\(Q_n\) lies in
\([-\vartheta_n-\omega_n,\vartheta_n+\omega_n]\cup[1-r-\omega_n,1+r+\omega_n]\);
its limit points lie in \(\{0\}\cup[1-r,1+r]\). (iii) Immediate from (i) and
Corollary \ref{cor:master-quantifier}.
\end{proof}

\begin{theorem}[strict symbol positivity \(\iff\) finite \emph{visible} count]
\label{thm:strict}
Strict symbol positivity at calendar \(a\)
holds iff \(\#\{n:J_n^{\mathrm{vis}}\neq\emptyset\}<\infty\), i.e.\
\(m_a^{\mathrm{vis}}:=\#\{j:\theta_j>0\}<\infty\). \textup{(}With noise
\(\omega_n\to0\) the equivalence holds \emph{verbatim}, with no exception
clause.\textup{)} In particular strict positivity says \emph{nothing} about
\(m-m_a^{\mathrm{vis}}\), the sub-resolution count.
\end{theorem}

\begin{proof}
(\(\Leftarrow\)) If \(J_n^{\mathrm{vis}}=\emptyset\) for \(n\geq N\),
then \(\spec(Q_n)\subseteq[1-r-\omega_n,1+r+\omega_n]\) for \(n\geq N\), so
\(\Lambda(Q)\subseteq[1-r,1+r]\) with gap \(c=1-r>0\) (Corollary
\ref{cor:master-quantifier}). (\(\Rightarrow\)) If infinitely many windows are
populated, each carries an eigenvalue
\(\mu_n\in[-\vartheta_n-\omega_n,-\vartheta_n+\omega_n]\) (Lemma
\ref{lem:window-spec}). The sequence \((\mu_n)\) over the populated windows is
bounded; any accumulation point \(\ell\) satisfies
\(\ell\leq\limsup_n(-\vartheta_n+\omega_n)\leq\limsup_n\omega_n=0\), and
\(\ell\in\Lambda(Q)\) by Lemma \ref{lem:symbol-spectrum}: strict positivity
fails, \emph{regardless of the comparison between \(\vartheta_n\) and
\(\omega_n\)}.
\end{proof}

\begin{remark}[a corrected exception clause]\label{rem:strict-corrected}
An earlier version of this theorem (and its source, an earlier analysis,
Thm.~5.3(ii)) allowed an exception for windows with
\(\vartheta_n\leq\omega_n\) eventually. That clause is \emph{false}
(counterexample: \(\omega_n=1/n\), one population in each even window with
\(\vartheta_n=1/(2n)\); the perturbed block eigenvalues lie in
\([-3/(2n),1/(2n)]\to0\), so \(0\in\spec(\sigma(Q))\) and there is no gap). The
corrected statement above is \emph{stronger}: sub-noise visible windows are
\emph{not} compatible with strict symbol positivity. If one wishes to discard
populations below the noise floor, the place to do it is the definition of
\(J_n^{\mathrm{vis}}\) (a threshold \(\theta_j>\omega\)), not an exception in
the theorem. Corollary \ref{cor:gap-quant}, Conjecture \ref{conj:lp134} and
Theorem \ref{thm:synergy} used only the regime \(\theta\geq c_0\) and are
unaffected.
\end{remark}

\begin{corollary}[the gap controls the visible count quantitatively]
\label{cor:gap-quant}
In the clean case, if \(\spec(Q_n)\subseteq[c,\infty)\) for \(n\geq N_0\), then
\(m_a^{\mathrm{vis}}\leq\sum_{n<N_0}|J_n^{\mathrm{vis}}|\), and with the
zero-density bound per window \(|J_n^{\mathrm{vis}}|\ll\Delta\log n\)
\cite[Thm.\ 5.8]{IwaniecKowalski}, \(m_a^{\mathrm{vis}}\ll N_0\log N_0\).
\end{corollary}

\begin{proof}
By Theorem \ref{thm:strict} all visible quadruples sit in windows \(n<N_0\); sum
the per-window count.
\end{proof}

\begin{theorem}[sub-resolution \(\Rightarrow\) the symbol is blind]
\label{thm:blind}
Let \(Z\) have \(m=\infty\) but \(\delta_j\,a(\gamma_j)\to0\) (all
sub-resolution, with \(J_n^{\mathrm{vis}}=\emptyset\)). Then, as sections over
the same field,
\[
Q-\Qaut\in\Jc,\qquad\text{hence}\qquad\sigma(Q)=\sigma(\Qaut)\ \text{in }\Cor,
\]
and if \(\Qaut\) has a gap (\(\spec\subseteq[1-r,1+r]\) modulo \(\omega\)), then
\(\sigma(Q)\) has \emph{strict} positivity -- while the defect of the
configuration, as the full form, is \(\delta=2m=\infty\) by Lemma
\ref{lem:mono-add}\textup{(ii)} applied to the hyperbolic blocks
\textup{(}which exist in the unrestricted test class but require support
exceeding the calendar\textup{)}. Strict symbol positivity does \emph{not} imply
\(m<\infty\).
\end{theorem}

\begin{proof}
By (V1) in perturbation form (Remark \ref{rem:theta}), the window difference
satisfies
\(\|Q_n-\Qaut_n\|\leq\sum_{j:|\gamma_j-n|\leq\Delta}\varphi(\delta_j a(\gamma_j))
+2\omega_n\) with \(\varphi(t)\to0\) as \(t\to0\). Since \(\gamma_j\to\infty\),
finitely many \(j\) per window, and \(\delta_j a(\gamma_j)\to0\), the per-window
sum tends to \(0\): the difference section is in \(\Jc\). The rest is Lemma
\ref{lem:symbol-norm} and (S1). The full-form defect is the sum of the
hyperbolic planes (Lemma \ref{lem:mono-add}(ii), each contributing a negative
direction in the unrestricted class).
\end{proof}

\begin{theorem}[calendar no-go: \(m<\infty\) is a boundary event over the
resolution filtration]\label{thm:nogo}
\begin{enumerate}
\item[\textup{(a)}] For \emph{every} calendar \(a\) there is a configuration
\(Z\) with \(m=\infty\) and \(\delta_j\,a(\gamma_j)\to0\) (take
\(\gamma_j:=j\), \(\delta_j:=\min(\tfrac14,1/(j\,a(\gamma_j)))\)): invisible at
\(a\) by Theorem \ref{thm:blind}.
\item[\textup{(b)}] For every configuration and quadruple, some calendar sees it
(\(a(\gamma_j)\geq1/\delta_j\)). Hence with
\(m_a:=\#\{j:\delta_j a(\gamma_j)\geq1\}\), \(\;m=\sup_a m_a\).
\item[\textup{(c)}] Consequently \(m<\infty\iff\exists M\,\forall a:m_a\leq M\),
a \(\Sigma_2\) statement over the filtration of resolutions, and \emph{no single
Weil--Toeplitz symbol (no fixed calendar) certifies it}.
\end{enumerate}
\end{theorem}

\begin{proof}
(a) \(\delta_j a(\gamma_j)\leq1/j\to0\); Theorem \ref{thm:blind}. (b)
Monotonicity of \(a\mapsto m_a\) and the pointwise observation. (c) Rewriting of
(b) together with (a).
\end{proof}

\begin{remark}[the trichotomy]\label{rem:trichotomy}
The dichotomy ``symbol or ideal'' is false; it is a trichotomy stratified by
\(\delta_j a(\gamma_j)\):
\[
\begin{array}{lll}
\text{fat }(\delta_j a(\gamma_j)\geq1\text{ i.o.}) &
\text{in the symbol: mass }\leq-c/2 & \text{Thm.\ \ref{thm:fat}}\\
\text{thin visible }(\theta_j\to0^+) &
\text{at the boundary }0\in\spec\sigma(Q) & \text{Thm.\ \ref{thm:thin}}\\
\text{sub-resolution }(\delta_j a(\gamma_j)\to0) &
\text{in the ideal: }\sigma(Q)=\sigma(\Qaut) & \text{Thm.\ \ref{thm:blind}}
\end{array}
\]
The strata are relative to the calendar (Theorem \ref{thm:nogo}); within
\(\mathfrak W\), it is the \emph{intermediate (thin-visible) stratum} that lives
on the symbol/compact boundary, accumulating at \(0\) -- the unique point
separating weak from strict positivity, which by Theorem \ref{thm:exact} is the
only point on which robust finiteness of the defect depends -- while the fat
stratum sits in the negative interior of the symbol and the sub-resolution
stratum in the ideal. The slogan ``the zeros live on the boundary'' is therefore
a theorem-in-model for the intermediate stratum only, not a global statement
about the zeros of \(\zeta\). The master quantifier acquires an
exact \(C^*\)-algebraic direction: it is the point \(0\in\spec_{\Cor}\).
\end{remark}

\subsection{Synergy of the two axes}
The Fredholm axis (symbol positivity) and the replication axis combine cleanly.

\begin{theorem}[replication \(\Rightarrow\) the symbol sees]\label{thm:synergy}
Suppose \(Z\) satisfies the quantitative replication axiom \emph{(R)}: if
\(m\geq1\), there is a quadruple \((\delta_0,\gamma_0)\in Z\) and heights
\(\tau_k\to\infty\) such that \(Z\) contains quadruples
\((\delta_k',\gamma_0+\tau_k)\) with \(\delta_k'\geq\delta_0/2\). Let \(a\) be any
calendar with \(a(\gamma)\to\infty\). Then \(m\geq1\) forces infinitely many fat
quadruples, hence
\(\spec(\sigma(Q))\cap(-\infty,-c_0/2]\neq\emptyset\). Contrapositively,
\[
\text{weak symbol positivity }\wedge\text{ (R)}\ \Longrightarrow\ m=0
\ \Longrightarrow\ \mathrm{RH}.
\]
\end{theorem}

\begin{proof}
The replicas have
\(\delta_k'a(\gamma_0+\tau_k)\geq(\delta_0/2)a(\gamma_0+\tau_k)\to\infty\), hence
are eventually supra-resolution with visibility \(\geq c_0\) by (V2). Theorem
\ref{thm:fat} applies. The contrapositive uses only weak positivity since
Theorem \ref{thm:fat} violates even weak positivity.
\end{proof}

The thin/sub-resolution stratum that kills the pure Fredholm route (Theorems
\ref{thm:blind}, \ref{thm:nogo}) is exactly the stratum that replication does
not produce: replicas inherit a fixed \(\delta_0\), and every unbounded
calendar eventually sees them. We record the missing piece this leaves open.

\begin{conjecture}[LP-134: resolution repulsion; new, non-RH, falsifiable]
\label{conj:lp134}
There exist a calendar \(a\) (e.g.\ \(a(\gamma)=C\log\gamma\)) and \(c>0\) such
that every nontrivial zero \(\rho=\half+\delta+i\gamma\) of \(\zeta\) with
\(\delta\neq0\) satisfies \(|\delta|\geq c/a(\gamma)\). \emph{(``Zeros cannot be
asymptotically sub-resolution.'')}
\end{conjecture}

Under LP-\ref{conj:lp134} the thin and sub-resolution strata are empty, and by
Theorems \ref{thm:fat}, \ref{thm:strict} the Fredholm half becomes a single
\(C^*\)-statement: \(m<\infty\iff\) strict symbol positivity at one calendar.
LP-\ref{conj:lp134} is \emph{not} RH (a quadruple at \(\delta=\tfrac14\)
satisfies it and violates RH), follows from nothing known (classical zero-free
regions repel from \(\operatorname{Re}s=1\), not from \(\operatorname{Re}s=\half\);
density theorems bound how many, not how close), and is of the right shape -- a
rigidity dichotomy.

\section{The localized signature and the logistic law}\label{sec:logistic}

We now develop the second, independent contribution: an exact logistic law for
the tensor tower of window indices, and an RH criterion in its localized
form (one branch a theorem, the other conditional on the Paley--Wiener uniformity problem). We work in
non-standard coordinates, using a non-principal ultrafilter
\(\mathcal U\) on \(\N\), the ultrapower \({}^*\R=\R^\N/\mathcal U\), and the
standard part \(\sst\) on the limited elements \cite{Robinson}.

\subsection{The logistic law}

\begin{theorem}[exact logistic law of the tensor tower]\label{thm:logistic}
Let \(\bar Q\) be a nondegenerate hermitian form of signature \((p,q)\) on a
space of dimension \(\bar n=p+q\), and set
\(x_k:=\negind(\bar Q^{\otimes 2^k})/\bar n^{2^k}\). Then for all \(k\geq0\), as
an identity of rationals,
\[
x_{k+1}=2x_k(1-x_k),\qquad 1-2x_k=(1-2x_0)^{2^k}.
\]
Consequently, for any scheme of forms with \(\sigma^0:=\sst([x_0(N)])\),
\[
\sigma^{(k)}:=\sst([x_k(N)])=\tfrac12\bigl(1-(1-2\sigma^0)^{2^k}\bigr),
\qquad
\lim_{k\to\infty}\sigma^{(k)}=
\begin{cases}0,&\sigma^0\in\{0,1\},\\[2pt]\tfrac12,&\sigma^0\in(0,1).\end{cases}
\]
\end{theorem}

\begin{proof}
By Lemma \ref{lem:whitney} applied to
\(\bar Q^{\otimes 2^{k+1}}=(\bar Q^{\otimes 2^k})\otimes(\bar Q^{\otimes 2^k})\),
with signature \((P_k,K_k)\), \(P_k+K_k=\bar n^{2^k}\),
\[
K_{k+1}=2P_kK_k,\qquad P_{k+1}=P_k^2+K_k^2.
\]
Dividing by \(\bar n^{2^{k+1}}=(\bar n^{2^k})^2\) gives
\(x_{k+1}=2(P_k/\bar n^{2^k})(K_k/\bar n^{2^k})=2x_k(1-x_k)\). Then
\(1-2x_{k+1}=1-4x_k+4x_k^2=(1-2x_k)^2\), so \(1-2x_k=(1-2x_0)^{2^k}\). Each
\(x_k(N)\in[0,1]\) is limited; \(\sst\) is a ring homomorphism on limited
elements, so the polynomial identity \(1-2x_{k+1}=(1-2x_k)^2\) passes to
standard parts for each fixed \(k\) (\(k\) standard; the identity is used
finitely often). Limits: if \(\sigma^0\in(0,1)\) then \(|1-2\sigma^0|<1\) and
\((1-2\sigma^0)^{2^k}\to0\); if \(\sigma^0=0\), \(\sigma^{(k)}\equiv0\); if
\(\sigma^0=1\), \(1-2\sigma^0=-1\) and \((-1)^{2^k}=1\) for \(k\geq1\), so
\(\sigma^{(k)}=0\) -- consistent, since the tensor square of a negative-definite
form is positive-definite.
\end{proof}

\begin{remark}[the global signature collapses]\label{rem:collapse}
The naive global signature -- normalizing the window index by the full window
dimension \(n_N\to\infty\) -- collapses to \(\sigma^0=0\) under \(\neg\)RH of
finite defect, because capturing an off\(\times\)on negative plane costs a full
dimensional factor \((w_\sigma\log T)(w_\tau\log T)\) while the off slot is rigid
(\(\leq4\bar m\) vectors): \(\kappa/n\lesssim\bar m/\max(w\log T)\to0\). The same
tensor identity that amplifies the tower (Theorem \ref{thm:logistic}) dilutes the
density of an individual defect. We call this \emph{density invisibility}: a
defect of finite cardinal is null for every asymptotic-density functional. This
forces the \emph{localized} construction below, where the per-cell dimension is
bounded and does not dilute. (The collapse is proved in the analysis preceding this paper, not in \cite{P40}; we do not re-prove it, and we
do not use the global signature as a criterion.)
\end{remark}

\subsection{The localized signature}

\begin{definition}[micro-resolution cell]\label{def:cell}
Fix a \emph{profile} \(p=(D,w,M,b,c,c')\) with \(D\in\N\) (degree), \(w>0\)
(width), \(M\in\N\) (decay order), \(b\in\mathbb Q_{>0}\) (caliber),
\((c,c')\in\mathbb Q^2\) (centers). The \emph{cell} is
\[
V_p:=\operatorname{span}\bigl\{h_{g_a,\varphi_b}:0\leq a,b\leq D\bigr\},\quad
\hat g_a(\sigma):=B^{(M)}_w(\sigma-ic)\sigma^a,\quad
\hat\varphi_b(\tau):=B^{(M)}_w(\tau-ic')\tau^b,
\]
where \(B^{(M)}_w\) is the Mellin transform of a fixed even bump
\(g_0^{(M,w)}\in C_c^\infty(\R_{>0})\) with
\[
|B^{(M)}_w(z)|\leq C_M e^{a_0|\operatorname{Re}z-1/2|}
(1+|\operatorname{Im}z|/w)^{-M},\qquad B^{(M)}_w(0)=1
\tag{$\ast$}
\]
(Paley--Wiener; \(a_0=a_0(M,w)\) the fixed log-support). Then \(\dim V_p\leq
(D+1)^2\), and \(V_p\) is test-accessible (the profile mentions no zeros). The
\emph{filled cell} is \(\tilde V_p:=V_p\oplus\C^{(D+1)^2}\) with
\(\tilde Q_p:=(Q_2|_{V_p})\oplus I\) (positive filler of dimension equal to the
cell's cap).
\end{definition}

\begin{definition}[localized scheme]\label{def:loc-scheme}
For each \(N\), the \emph{grid} \(\mathcal G_N\) is the finite computable set of
profiles \(p\) with \(D,M\leq N\), \(w,b\in\{2^{-j}:j\leq N\}\), and centers in
\((b\Z)^2\cap[-2^N,2^N]^2\). Set
\[
x_0(N):=\max_{p\in\mathcal G_N}\frac{\negind(\tilde Q_p)}{\dim\tilde V_p}
\in\mathbb Q\cap[0,\tfrac12],
\]
and \(\sigma_{\mathrm{loc}}:=\sst([x_0(N)])\). For each \(N\) fix a winning
profile \(p^*_N\) (maximizing the ratio), let \(\bar Q_N\) be the reduced
nondegenerate form of \(\tilde Q_{p^*_N}\) (quotient by the radical) and
\(\bar n_N:=\dim\tilde V_{p^*_N}-\dim\operatorname{rad}(\tilde Q_{p^*_N})\) its
dimension, and set
\[
\begin{gathered}
x_{\mathrm{base}}(N):=\frac{\negind(\bar Q_N)}{\bar n_N},\qquad
x_k(N):=\frac{\negind(\bar Q_N^{\otimes2^k})}{\bar n_N^{2^k}},\\[2pt]
\sigma_{\mathrm{base}}^{\,\mathcal U}:=\sst([x_{\mathrm{base}}(N)]),\qquad
\sigma_{\mathrm{loc}}^{(k)}:=\sst([x_k(N)]).
\end{gathered}
\]
Note that \(x_{\mathrm{base}}(N)\neq x_0(N)\) in general when the radical is
nontrivial (the denominators differ), and that
\(\sigma_{\mathrm{base}}^{\,\mathcal U}\) and the intermediate floors
\(\sigma_{\mathrm{loc}}^{(k)}\), \(k\geq1\), are not claimed to be independent
of \(\mathcal U\): only the limit dichotomy of Proposition
\ref{prop:loc-logistic} is.
\end{definition}

\begin{proposition}[well-definedness and autonomy]\label{prop:loc-welldef}
\textup{(a)} \(x_0(N)\) is determined by the arithmetic data without zero
positions: every Gram entry is \(Q_2(h,h')\) (autonomy, Hypothesis
\ref{hyp:autonomy}), a computable real number (a prime sum), and the grid is
finite. Since the entries are computable reals rather than rationals,
\(\negind\) is \emph{not} claimed to be computed exactly; what holds, and what
Section \ref{sec:honest} uses, is that the predicate
\(\negind(\tilde Q_p)\geq1\) is \emph{semidecidable} by interval arithmetic: a
strictly negative eigenvalue is certified at finite precision. \textup{(b)}
\(\sigma_{\mathrm{loc}}\in[0,\tfrac12]\) exists for every \(\mathcal U\).
\textup{(c)} \(x_0(N)\) is nondecreasing in \(N\) (\(\mathcal G_N\subseteq
\mathcal G_{N+1}\)), so it converges and \(\sigma_{\mathrm{loc}}=\lim_N x_0(N)\)
is independent of \(\mathcal U\).
\end{proposition}

\begin{proof}
(a) The class contains \(B_w\cdot(\text{monomials})\) by closure under Mellin
multiplication \cite{P40}; the filler is an explicit block; a strictly negative
eigenvalue of a hermitian matrix with computable-real entries is certified at
finite precision by continuity of eigenvalues (Theorem \ref{thm:weyl-stab}).
(b) \(\sst\) exists on limited elements. (c) Monotone
bounded; \(\sst\) of a convergent sequence is its limit, for every
\(\mathcal U\).
\end{proof}

\begin{proposition}[the logistic law, localized and repaired]
\label{prop:loc-logistic}
For every \(\mathcal U\) and standard \(k\geq0\),
\[
\sigma_{\mathrm{loc}}^{(k)}
=\half\bigl(1-(1-2\sigma_{\mathrm{base}}^{\,\mathcal U})^{2^k}\bigr),
\qquad
\sigma_{\mathrm{base}}^{\,\mathcal U}\in
[\sigma_{\mathrm{loc}},\,2\sigma_{\mathrm{loc}}]\cap[0,\tfrac12].
\]
In particular, \emph{independently of \(\mathcal U\)}:
\(\lim_k\sigma_{\mathrm{loc}}^{(k)}=0\) if \(\sigma_{\mathrm{loc}}=0\), and
\(=\half\) if \(\sigma_{\mathrm{loc}}\in(0,\tfrac12]\).
\end{proposition}

\begin{proof}
For each \(N\), \(\bar Q_N\) is finite-dimensional nondegenerate; Theorem
\ref{thm:logistic} applies coordinatewise with base \(x_{\mathrm{base}}(N)\),
and \(\sst\) is a ring homomorphism on limited elements, giving the closed
formula in \(\sigma_{\mathrm{base}}^{\,\mathcal U}\). The bracketing: the filler
\(I_{(D+1)^2}\) is nondegenerate and orthogonal to the compressed block, so
\(\operatorname{rad}(\tilde Q_p)=\operatorname{rad}(Q_2|_{V_p})\oplus0\) and
\(\bar n_N\geq(D+1)^2=\tfrac12\dim\tilde V_{p^*_N}\), whence
\(\dim\tilde V_{p^*_N}/\bar n_N\in[1,2]\); since
\(\negind(\bar Q_N)=\negind(\tilde Q_{p^*_N})\) (the radical carries no negative
direction) and \(\sst\) is multiplicative on limited elements,
\(\sigma_{\mathrm{loc}}\leq\sigma_{\mathrm{base}}^{\,\mathcal U}
\leq2\sigma_{\mathrm{loc}}\). Moreover
\(x_{\mathrm{base}}(N)\leq(\bar n_N-(D+1)^2)/\bar n_N\leq\half\), so
\(\sigma_{\mathrm{base}}^{\,\mathcal U}\in[0,\half]\) and the case
\(\sigma^0=1\) (the only one breaking the dichotomy in Theorem
\ref{thm:logistic}) is excluded. The dichotomy: \(\sigma_{\mathrm{loc}}=0\iff
\sigma_{\mathrm{base}}^{\,\mathcal U}=0\) (bracketing), in which case the tower
is identically \(0\); and \(\sigma_{\mathrm{loc}}>0\Rightarrow
\sigma_{\mathrm{base}}^{\,\mathcal U}\in(0,\half]\Rightarrow
|1-2\sigma_{\mathrm{base}}^{\,\mathcal U}|<1\Rightarrow
\lim_k\sigma_{\mathrm{loc}}^{(k)}=\half\), for every \(\mathcal U\).
\end{proof}

\begin{remark}[corrected]\label{rem:radical-bug}
An earlier version of this proposition asserted the closed formula with
\(\sigma_{\mathrm{loc}}\) itself as the base of the tower. That is \emph{false}
as stated whenever the winning Gram is degenerate: the definition then carried
two incompatible normalizations of \(x_0(N)\) (denominator
\(\dim\tilde V_p\) versus the reduced dimension \(\bar n_N\)). The repair above
(found in an adversarial audit) keeps the \(0\)-versus-\(\half\)
dichotomy and the \(\mathcal U\)-independence of the \emph{limit}; the closed
formula holds in terms of \(\sigma_{\mathrm{base}}^{\,\mathcal U}\), and the
intermediate floors are not claimed canonical.
\end{remark}

\subsection{Stability of the cell index}
These three perturbation lemmas make a computable mesh suffice to capture a plane
whose exact position is unknown.

\begin{theorem}[Weyl stability]\label{thm:weyl-stab}
For hermitian \(A,E\) on a finite-dimensional space, with eigenvalues ordered
\(\lambda_1\geq\dots\geq\lambda_n\), \(|\lambda_k(A+E)-\lambda_k(A)|\leq\|E\|\)
for all \(k\). In particular, if \(A\) has \(\kappa\) eigenvalues \(\leq-\mu\)
and \(\|E\|<\mu\), then \(\negind(A+E)\geq\kappa\).
\end{theorem}

\begin{proof}
Courant--Fischer: \(\lambda_k(A+E)=\min_{\operatorname{codim}S=k-1}
\max_{0\neq v\in S}R_{A+E}(v)\) with \(R\) the Rayleigh quotient; for all \(v\),
\(|R_{A+E}(v)-R_A(v)|\leq\|E\|\). Take the optimal \(S\) for \(A\) and swap
\(A\leftrightarrow A+E\) \cite{Bhatia}.
\end{proof}

\begin{theorem}[bounded-rank stability]\label{thm:rank-stab}
\(|\negind(A+E)-\negind(A)|\leq\rank E\).
\end{theorem}

\begin{proof}
Let \(N\) be negative for \(A\) with \(\dim N=\negind(A)\). On
\(N\cap\ker E\) (codimension \(\leq\rank E\) in \(N\)),
\((A+E)|_{N\cap\ker E}=A|_{N\cap\ker E}<0\), so
\(\negind(A+E)\geq\negind(A)-\rank E\). Swap \(A\leftrightarrow A+E\) and
\(E\leftrightarrow-E\).
\end{proof}

\begin{theorem}[Lipschitz Gram under cell shift]\label{thm:lipschitz}
For a profile \(p\), let \(G(\vartheta)\) be the Gram of \(Q_2\) on the cell
with centers \((c+\vartheta_1,c'+\vartheta_2)\). Then \(G\) is differentiable in
\(\vartheta\) with \(\|\partial_{\vartheta_i}G(\vartheta)\|\leq L_p<\infty\),
\(L_p\) computable from the arithmetic data.
\end{theorem}

\begin{proof}
Center shift is the modulation \(g\mapsto g\cdot x^{i\vartheta}\), whose
\(\vartheta\)-derivative inserts \(i(\log x)\)g, a test of the same class
\cite{P40}. The spectral side \eqref{eq:speca} (below) converges absolutely with
local uniformity in \(\vartheta\) by the decay \((\ast)\) \cite{P40}, so one may
differentiate under the sum; the derivative inserts a linear factor, i.e.\ a
derived test of the class, whose values are finite and computable. The norm of a
\((D+1)^2\times(D+1)^2\) hermitian matrix is at most its dimension times the
maximal entry.
\end{proof}

\begin{theorem}[the aligned cell survives in the mesh]\label{thm:mesh}
Let \(p^*\) be an ideal profile (real centers, not necessarily in the mesh) with
\(G(p^*)\) having \(\kappa\) eigenvalues \(\leq-\mu<0\). Let \(p\in\mathcal G_N\)
be the nearest mesh profile, \(|\mathrm{centers}(p)-\mathrm{centers}(p^*)|\leq
b\sqrt2\). If \(b\leq\mu/(4\sqrt2\,L_{p^*})\), then \(\negind(G(p))\geq\kappa\).
In particular, for every ideal profile with \((\mu,L)\) fixed there is an
\(N_0\), explicit in \((\mu,L,\mathrm{position})\), such that \(\mathcal G_N\)
contains a cell of index \(\geq\kappa\) for all \(N\geq N_0\). \textup{(}We say
``explicit'', not ``computable'': the inputs \((\mu,L,\mathrm{position})\) are
not effectively knowable; cf.\ Section \ref{sec:honest}\textup{(a)}.\textup{)}
\end{theorem}

\begin{proof}
By Theorem \ref{thm:lipschitz} and the mean value theorem,
\(\|G(p)-G(p^*)\|\leq L_{p^*}\cdot b\sqrt2\cdot2=2\sqrt2\,L_{p^*}b\leq\mu/2<\mu\);
apply Theorem \ref{thm:weyl-stab}. Existence of \(N_0\): the mesh has caliber
\(2^{-N}\to0\) and reach \(2^N\to\infty\). \textup{(}The earlier bound
\(b\leq\mu/(2\sqrt2\,L)\) gave only \(\leq\mu\), not the strict inequality
Theorem \ref{thm:weyl-stab} requires; corrected.\textup{)}
\end{proof}

\subsection{Realizability under \texorpdfstring{$\neg$}{not-}RH}
We now import the spectral side and prove the realizability theorem. From
\cite{P40,P35} we have, on the span of the (doubled) test class, the spectral
side
\begin{equation}\label{eq:speca}
Q_2(h,h')=\sum_{(\rho_1,\rho_2)}\hat h(\rho_1,\rho_2)\,
\overline{\hat h'(\iota(\rho_1,\rho_2))},\quad
\iota(\rho_1,\rho_2):=(1-\bar\rho_1,1-\bar\rho_2),
\end{equation}
absolutely convergent, with \(\hat h_{g,\varphi}(\rho_1,\rho_2)=
\half\hat g(\sigma)\hat\varphi(\tau)\), \(\sigma=(\rho_1+\rho_2)/2\),
\(\tau=(\rho_1-\rho_2)/2\); pairs fixed by \(\iota\) (both online) contribute
\(|\hat h|^2\geq0\), and free orbits \(\{P,\iota P\}\) contribute hyperbolic
planes of signature \((1,1)\); under RH the sum is diagonal and \(Q_2\geq0\)
\cite{P40,P35}. We use \(N(t+1)-N(t)\ll\log(2+t)\) \cite[Thm.\
5.8]{IwaniecKowalski}.

\begin{lemma}[tail of the concentrated test]\label{lem:tail}
Let \(h\in V_p\) with \(|\hat g(\sigma)|\leq C_g(1+|\operatorname{Im}\sigma-
\gamma_0|)^{-M+D}\) and \(|\hat\varphi(\tau)|\leq C_\varphi(1+|\operatorname{Im}
\tau|)^{-M+D}\) on \(|\operatorname{Re}|\leq1\). Let \(\Omega_r\) be the pairs
with \(\max(|\operatorname{Im}\sigma-\gamma_0|,|\operatorname{Im}\tau|)\geq r\).
If \(M-D\geq4\), then
\[
\sum_{(\rho_1,\rho_2)\in\Omega_r}|\hat h(\rho_1,\rho_2)|\,
|\hat h(\iota(\rho_1,\rho_2))|\leq
C_g^2C_\varphi^2\,\frac{C^*(M-D)\,(\log(2+|\gamma_0|))^2}{(1+r)^{M-D-4}},
\]
\(C^*\) absolute. \emph{Warning (corrected): the constants
\(C_g,C_\varphi\) depend on \(M\) through the bump \((\ast)\); the right-hand
side tends to \(0\) as \(M\to\infty\) at fixed \(r\) only if the family of
bumps keeps \(C_g(M)^2C_\varphi(M)^2(1+r)^{-(M-D-4)}\to0\) -- this uniformity is
exactly Hypothesis \ref{hyp:pw} below and is \emph{not} proved here.}
\end{lemma}

\begin{proof}
Write \(u=\operatorname{Im}\sigma=(\gamma_1+\gamma_2)/2\),
\(v=\operatorname{Im}\tau=(\gamma_1-\gamma_2)/2\), so \(\gamma_1=u+v\),
\(\gamma_2=u-v\). With \(a=|\gamma_1-\gamma_0|\), \(b=|\gamma_2-\gamma_0|\), the
elementary inequality \((1+\tfrac{a+b}2)^2\geq(1+a)(1+b)\) (since
\((1+\tfrac{a+b}2)^2-(1+a)(1+b)=(\tfrac{a-b}2)^2\geq0\)) and
\(|u-\gamma_0|+|v|\geq(a+b)/2\) give
\[
(1+|u-\gamma_0|)(1+|v|)\geq 1+|u-\gamma_0|+|v|\geq1+\tfrac{a+b}2
\geq\sqrt{(1+a)(1+b)}.
\]
The point \(\iota(\rho_1,\rho_2)\) has the same ordinates
(\(\operatorname{Im}(1-\bar\rho)=\operatorname{Im}\rho\)), so with \(m:=M-D\),
\[
|\hat h(P)||\hat h(\iota P)|\leq\tfrac14 C_g^2C_\varphi^2
\bigl[(1+|u-\gamma_0|)(1+|v|)\bigr]^{-2m}
\leq\tfrac14 C_g^2C_\varphi^2(1+a)^{-m}(1+b)^{-m}.
\]
The double sum factors; the single sum, grouped over unit intervals
\(|\gamma-\gamma_0|\in[j,j+1)\) with \(N(t+1)-N(t)\ll\log(2+t)\), is
\(\sum_\rho(1+|\gamma-\gamma_0|)^{-m}\leq C'\zeta(m-1)\log(2+|\gamma_0|)\),
convergent for \(m\geq3\). On \(\Omega_r\), one product is \(\geq1+r\); extract
\((1+r)^{-(m-4)}\) and retain exponent \(2\) in each factor for convergence.
\end{proof}

\begin{assumption}[Paley--Wiener uniformity; the Paley--Wiener uniformity problem]\label{hyp:pw}
There is a family of even bumps \(\{g_0^{(M,w)}\}_M\) satisfying \((\ast)\)
such that, in the witness construction of Theorem \ref{thm:realiz} (fixed
configuration, fixed \(D\) and \(w\)), the full tail budget of Lemma
\ref{lem:tail} at \(r=1\) tends to \(0\) as \(M\to\infty\):
\[
C_1(M)^2\,C_g(M)^2\,C_\varphi(M)^2\;
C^*(M-D)\,(\log(2+|\gamma_0|))^2\;2^{-(M-D-4)}\;\longrightarrow\;0,
\]
where \(C_g(M),C_\varphi(M)\) absorb the Paley--Wiener constant \(C_M\) and the
strip factor \(e^{a_0(M,w)}\) of \((\ast)\), and \(C_1(M)\) is the Vandermonde
coefficient norm of Lemma \ref{lem:vandermonde} (which depends on \(M\) through
\(b_0(M)=\min_i|B^{(M)}(z_i-ic)|\)).
\end{assumption}

\begin{remark}[why Hypothesis \ref{hyp:pw} is a genuine open gap:
Paley--Wiener rigidity]\label{rem:pw-rigidity}
For any \emph{fixed} nonzero bump (fixed log-support \(a_0\)), the optimal
constants \(C_M\) in \((\ast)\) grow superexponentially: for every \(A>1\),
\(C_M>A^M\) for all large \(M\). Indeed, if \(C_M\leq A^M\) for infinitely many
\(M\), then for \(|t|>Aw\) one would have \(|B(t)|\leq(Aw/|t|)^M\to0\) along
those \(M\), forcing \(B\equiv0\) on \(\{|t|>Aw\}\); but \(B\) is entire
(Paley--Wiener, compact support) and not identically zero, so its real zeros
are discrete -- contradiction. Hence \(C_M^4\,2^{-M}\to\infty\): \emph{no fixed
bump satisfies Hypothesis \ref{hyp:pw}}, and the standard \(M\)-dependent
families (powers of \(\operatorname{sinc}\), support \(a_0(M)\) growing) fail
as well -- the strip factor \(e^{c\,a_0(M)}\) defeats the polynomial decay in
\(\operatorname{Im}\) at fixed distance \(r=1\). Closing the gap requires
either a family of bumps with the stated uniformity, or a different proof
mechanism (the classical exponential-dominance route of Weil's criterion, in
which the off-zero term grows like \(e^{2\delta_0 a_0}\) with \(a_0\to\infty\)
and absorbs the \(M\)-dependent constants). This is \emph{the Paley--Wiener uniformity problem}, open.
We emphasize the precise status: the \emph{statement} of Theorem
\ref{thm:realiz} is not refuted -- Lemma \ref{lem:tail} is only an upper bound,
and the classical mechanism makes the statement plausible -- but its proof here
is conditional on Hypothesis \ref{hyp:pw}.
\end{remark}

\begin{lemma}[Vandermonde with bump]\label{lem:vandermonde}
Let \(z_1,\dots,z_q\in\C\) be distinct and \(B\) entire with \(B(z_i)\neq0\). If
\(D\geq q-1\), the evaluations \(f\mapsto f(z_i)\) are linearly independent on
\(\operatorname{span}\{B(\sigma)\sigma^a:0\leq a\leq D\}\); for any data
\((y_1,\dots,y_q)\) there is an element with \(f(z_i)=y_i\) and coefficient norm
\(\leq C(z,\min_i|B(z_i)|)\|y\|\). Moreover, for
\(B=B^{(M)}_w(\cdot-ic)\not\equiv0\), the centers \(c\) with \(B(z_i-ic)=0\) for
some \(i\) are discrete; every open interval of centers contains a \(c\) with
\(\min_i|B(z_i-ic)|\geq b_0>0\), stable on a ball.
\end{lemma}

\begin{proof}
The matrix \((B(z_i)z_i^a)_{i,a}=\operatorname{diag}(B(z_i))\cdot(z_i^a)\) has
rank \(q\) (Vandermonde with distinct nodes, invertible diagonal). The function
\(c\mapsto B(z_i-ic)\) is entire and not identically zero, so its real zeros are
discrete; a finite union of discrete sets is discrete; \(B\) is continuous and
nonzero on a compact set avoiding the zeros.
\end{proof}

\begin{theorem}[realizability under \(\neg\)RH, conditional on the Paley--Wiener uniformity problem]
\label{thm:realiz}
Assume \(\neg\)RH \emph{and Hypothesis \ref{hyp:pw}}. Let
\(\rho^0=\half+\delta_0+i\gamma_0\) be the off-critical
zero of minimal height, multiplicity \(m_0\geq1\). Then there are a standard
dimension \(D_0\) and a standard \(N_0\) such that for all \(N\geq N_0\),
\(x_0(N)\geq 1/(2D_0)>0\), and hence for every non-principal \(\mathcal U\),
\[
\sigma_{\mathrm{loc}}\geq\frac{1}{2D_0}>0,\qquad
\lim_{k\to\infty}\sigma_{\mathrm{loc}}^{(k)}=\tfrac12.
\]
\end{theorem}

\begin{proof}
\emph{Step 0 (target plane).} Take \(P_0=(\rho^0,\rho^0)\) with coordinates
\((\sigma_0,\tau_0)=(\half+\delta_0+i\gamma_0,0)\) and
\(\iota P_0=(1-\bar\rho^0,1-\bar\rho^0)\) with
\((\sigma_0',0)=(\half-\delta_0+i\gamma_0,0)\). Since \(\delta_0\neq0\),
\(\iota P_0\neq P_0\): a free orbit, contributing
\(2m_0^2\operatorname{Re}[\hat h(P_0)\overline{\hat h(\iota P_0)}]\) to
\eqref{eq:speca}.

\emph{Step 1 (the sets to annihilate).} Let \(S\) be the finite set of zeros
\(\rho\) with \(||\gamma|-|\gamma_0||\leq2\), and \(Z_\sigma\) the finite set of
\(\sigma=(\rho_1+\rho_2)/2\) over pairs with both members in \(S\), excluding
\(\sigma_0,\sigma_0'\). Let \(Z_\tau\) be the finite set of
\(\tau=(\rho_1-\rho_2)/2\) over the \emph{resonant pairs}: pairs with both
members in \(S\) and \(\sigma\in\{\sigma_0,\sigma_0'\}\), other than the target
orbit \(\{P_0,\iota P_0\}\). Every resonant pair has \(\tau\neq0\) (\(\tau=0\)
together with \(\sigma=\sigma_0\) forces \(\rho_1=\rho_2=\rho^0\), the target).
Put \(D:=\max(|Z_\sigma|+1,\,|Z_\tau|)\); all standard, fixed by the
configuration. \textup{(}The annihilation of the resonant pairs is a correction:
an earlier version of this proof omitted them, and under general
\(\neg\)RH configurations they contribute terms of uncontrolled
sign.\textup{)}

\emph{Step 2 (the cell and witness).} Choose the ideal profile \(p^*\) with
degree \(D\), \(\sigma\)-center \(c^*\) where Lemma \ref{lem:vandermonde} gives
\(\min|B(z-ic^*)|\geq b_0\) over \(Z_\sigma\cup\{\sigma_0,\sigma_0'\}\), and
\(\tau\)-center \(c'^*\) where it gives the same over \(Z_\tau\cup\{0\}\). In
\(V_{p^*}\) build \(v=h_{g_v,\varphi_v}\) with
\(\hat g_v=0\) on \(Z_\sigma\), \(\hat g_v(\sigma_0)=1\),
\(\hat g_v(\sigma_0')=-1\), and \(\hat\varphi_v=0\) on \(Z_\tau\),
\(\hat\varphi_v(0)=1\) (Lemma \ref{lem:vandermonde} on each side; coefficient
norm \(\leq C_1\), standard at fixed \(M\) but \emph{\(M\)-dependent}, cf.\
Hypothesis \ref{hyp:pw}). The data \((1,-1)\) fixes the
phase: the target plane contributes
\(2m_0^2\operatorname{Re}[(\half\cdot1\cdot1)\overline{(\half\cdot(-1)\cdot1)}]
=-m_0^2/2<0\).

\emph{Step 3 (budget of \(Q_2(v,v)\)).} Decompose \eqref{eq:speca}: (a) the orbit
\(\{P_0,\iota P_0\}\) contributes \(-m_0^2/2\); (b) pairs with both members in
\(S\) other than the target: either \(\sigma\in Z_\sigma\), where
\(\hat g_v=0\), or \(\sigma\in\{\sigma_0,\sigma_0'\}\) (resonant), where
\(\tau\in Z_\tau\) and \(\hat\varphi_v=0\) -- in both cases the contribution is
\(0\) (the case split is exhaustive, and the \(\iota\)-image adds no new case:
\(\sigma(\iota P)=1-\bar\sigma\) swaps \(\sigma_0\leftrightarrow\sigma_0'\),
already covered); (c) the rest lie in \(\Omega_1\) about \((\gamma_0,0)\) (some
member outside \(S\)), and by Lemma \ref{lem:tail} with \(r=1\) the bound is
the budget of Hypothesis \ref{hyp:pw}, which \emph{by that hypothesis -- and
only by it (Remark \ref{rem:pw-rigidity}); this is GAP-A, the conditional point
of the theorem} -- tends to \(0\) as \(M\to\infty\). Choose \(M=M_0\) standard
with \(|(c)|\leq m_0^2/4\). Then
\[
Q_2(v,v)\leq-\frac{m_0^2}{2}+0+\frac{m_0^2}{4}=-\frac{m_0^2}{4}<0,
\]
so \(\negind(Q_2|_{V_{p^*}})\geq1\), with a standard spectral margin \(-\mu_0\).

\emph{Step 4 (ideal cell to grid).} By Theorem \ref{thm:mesh} (with
\(\kappa=1\), \(\mu=\mu_0\), \(L=L_{p^*}\) finite by Theorem \ref{thm:lipschitz},
all standard) there is a standard \(N_0\) such that for \(N\geq N_0\),
\(\mathcal G_N\) contains a profile \(p\) with \(\negind(Q_2|_{V_p})\geq1\).

\emph{Step 5 (conclusion).} For \(N\geq N_0\),
\(x_0(N)\geq\negind(\tilde Q_p)/\dim\tilde V_p\geq1/(2(D+1)^2)=:1/(2D_0)\) with
\(D_0=(D+1)^2\) standard (filler doubles dimension; the negative subspace
persists).
By Proposition \ref{prop:loc-welldef}(c),
\(\sigma_{\mathrm{loc}}=\lim x_0(N)\geq1/(2D_0)\) for every \(\mathcal U\);
Proposition \ref{prop:loc-logistic} gives \(\sigma_{\mathrm{loc}}^{(k)}\to\half\).
\end{proof}

\begin{remark}[the conceptual key; what replaces the realizability hypothesis]
\label{rem:key}
The index \(\negind(Q_2|_{V_p})\) is a property of the \emph{compressed Gram}.
Test-accessibility restricts the \emph{specification} of \(V_p\) (no zero
positions); the \emph{witness} \(v\) inside \(V_p\) is an object of the proof,
free to use all positions. Earlier work \cite{P40} used polynomial factors with
roots at zero positions and so appeared to inherit a realizability caveat (the
``Hypothesis D'' of \cite{P40}); here the same factors are rebuilt as linear
combinations \emph{inside} the test-accessible span \(\{B\cdot\sigma^a\}\)
(Lemma \ref{lem:vandermonde}). Localizing to the bounded height of the defect
dissolves the \(T\to\infty\) regime of that caveat -- but, as the
adversarial audit established, it does not dissolve the caveat
entirely: it \emph{transmutes} it into the uniformity in \(M\) of the
Paley--Wiener constant budget, Hypothesis \ref{hyp:pw} (the Paley--Wiener uniformity problem), a
strictly narrower and purely analytic statement, open. \emph{We claim no
quantitative bound}: \(D_0,N_0\) are not effective in \(\bar m\) -- they depend
on the unknown local configuration. That non-effectivity is exactly what the
criterion does not need.
\end{remark}

\begin{proposition}[RH branch]\label{prop:RHbranch}
Under RH, \(x_0(N)=0\) for all \(N\), \(\sigma_{\mathrm{loc}}=0\), and
\(\sigma_{\mathrm{loc}}^{(k)}\equiv0\).
\end{proposition}

\begin{proof}
Under RH the spectral side \eqref{eq:speca} is diagonal, so \(Q_2\geq0\) on the
span \cite{P40,P35}; every compression of a positive form is positive; the
filler is positive; \(\negind=0\) in every cell.
\end{proof}

\begin{theorem}[the logistic criterion]\label{thm:main}
For every non-principal ultrafilter \(\mathcal U\):
\begin{enumerate}
\item[\textup{(i)}] \textup{(theorem)}
\(\mathrm{RH}\Longrightarrow\sigma_{\mathrm{loc}}=0\Longrightarrow
\lim_k\sigma_{\mathrm{loc}}^{(k)}=0\);
\item[\textup{(ii)}] \textup{(conditional on Hypothesis \ref{hyp:pw},
the Paley--Wiener uniformity problem)}
\(\neg\mathrm{RH}\Longrightarrow\sigma_{\mathrm{loc}}>0\Longrightarrow
\lim_k\sigma_{\mathrm{loc}}^{(k)}=\tfrac12\).
\end{enumerate}
Hence, \emph{conditionally on Hypothesis \ref{hyp:pw}},
\[
\mathrm{RH}\;\Longleftrightarrow\;\sigma_{\mathrm{loc}}=0\;\Longleftrightarrow\;
\lim_{k\to\infty}\sigma_{\mathrm{loc}}^{(k)}=0\;\Longleftrightarrow\;
\lim_{k}\sigma_{\mathrm{loc}}^{(k)}\neq\tfrac12 .
\]
Unconditionally, \(\lim_k\sigma_{\mathrm{loc}}^{(k)}\in\{0,\half\}\) with
\(0\iff\sigma_{\mathrm{loc}}=0\), independently of \(\mathcal U\).
\end{theorem}

\begin{proof}
(i) Proposition \ref{prop:RHbranch} and Proposition \ref{prop:loc-logistic}.
(ii) Theorem \ref{thm:realiz} and Proposition \ref{prop:loc-logistic}. The
\(\mathcal U\)-independence of \(\sigma_{\mathrm{loc}}\) is Proposition
\ref{prop:loc-welldef}(c); the unconditional dichotomy is Proposition
\ref{prop:loc-logistic}.
\end{proof}

\begin{remark}[\(\mathcal U\) is tamed for free]
The dichotomy of Theorem \ref{thm:main} is independent of the ultrafilter:
\(\sigma_{\mathrm{loc}}\) is a genuine limit (Proposition
\ref{prop:loc-welldef}(c)). The non-canonicity one might fear -- ``which
ultrafilter concentrates where the defect lives?'' -- is resolved by
monotonicity of the growing grid, not by a choice.
\end{remark}

\section{Honest scope}\label{sec:honest}

We state precisely what is and is not achieved, per the contract.

\subsection{(a) RH is not proved, and no progress toward a proof is claimed}
Theorem \ref{thm:main} is an \emph{equivalence conditional on Hypothesis
\ref{hyp:pw}} (its RH branch is an unconditional theorem), a reformulation of
RH. It does
not decide RH and provides no new route to deciding it. The standard part
\(\sigma_{\mathrm{loc}}=0\) is a \(\Pi_1\) statement over the sequence
\((x_0(N))_N\) -- ``for all large \(N\), \(x_0(N)\) small'' -- exactly as
non-effective as RH itself. The improvement is one of \emph{coordinates}: each
event \(x_0(N)\) is now finite linear algebra over arithmetic data (Proposition
\ref{prop:loc-welldef}(a)), so falsifying RH becomes ``certify, by interval
arithmetic, a negative eigenvalue
of an arithmetic Gram'' rather than ``integrate \(\zeta'/\zeta\) with controlled
error''. This is a second semidecision machine alongside the classical numerical
one -- with the caveat that its \emph{completeness} (that \(\neg\)RH forces the
machine to halt) is Theorem \ref{thm:realiz}, hence conditional on
the Paley--Wiener uniformity problem; RH gains nothing.

\subsection{(b) The computability of the symbol is the classical wall, moved}
The trichotomy (Section \ref{sec:trichotomy}) locates the obstruction at the
single point \(0\in\spec\sigma(Q)\), and the following identity shows that
deciding whether the symbol equals its autonomous prediction is the classical
wall in new coordinates.

\begin{theorem}[computing the symbol \(=\) uniformity over windows]
\label{thm:symbolcompute}
On the same field of test spaces,
\(\|\sigma(Q)-\sigma(\Qaut)\|_{\Cor}=\limsup_n\|Q_n-\Qaut_n\|\). Hence the
statement ``the symbol of the Weil form is the archimedean/near-prime
prediction'' is equivalent to \(\|Q_n-\Qaut_n\|\to0\): the second-order
deficit asymptotics, in operator norm per window, \emph{uniformly} over \(n\) on
the whole ball of tests.
\end{theorem}

\begin{proof}
Lemma \ref{lem:symbol-norm} applied to the difference section; the right-hand
side is the per-window deficit form against the window's tests
\cite{P40}.
\end{proof}

The right-hand side is the uniform second-order pair-correlation asymptotic --
the classical wall. What arithmetic computes unconditionally are the
\emph{Ces\`aro states} (the mean over windows) of \(\sigma(Q)\); strict
positivity is a statement of \emph{norm} (all states, including the pure states
supported on window subsequences aligned with the zeros). The gap between the
computable and the necessary is exactly the gap between the mean states and the
exceptional pure states of the corona. We record the abstract half of what is
provable.

\begin{proposition}[mean Ces\`aro positivity, abstract half]\label{prop:cesaro}
If, for every \(\varepsilon>0\),
\(N^{-1}\#\{n\leq N:\lambda_{\min}(Q_n)<-\varepsilon\}\to0\) (``almost every
window is \(\varepsilon\)-positive''), then for every positive test section
\(0\leq A_n\leq I\) and every Banach limit \(\mathcal L\),
\[
\mathcal L\text{-}\lim_N\frac1N\sum_{n\leq N}\frac{\tr(Q_nA_n)}{d_n}\geq0:
\]
all positive Ces\`aro states of \(\sigma(Q)\) are \(\geq0\).
\end{proposition}

\begin{proof}
If \(\lambda_{\min}(Q_n)\geq-\varepsilon\), then \(Q_n\geq-\varepsilon I\), so
\(\tr(Q_nA_n)\geq-\varepsilon\tr(A_n)\geq-\varepsilon d_n\) (using
\(\tr(CA)\geq0\) for \(C,A\geq0\)); otherwise the trivial bound
\(\tr(Q_nA_n)\geq-\sup_k\|Q_k\|\,d_n\). The bad windows have vanishing density,
so the Banach limit of a sequence \(\geq-\varepsilon-o(1)\) is
\(\geq-\varepsilon\); let \(\varepsilon\to0\).
\end{proof}

The hypothesis is the mean-value / large-sieve output (almost every window),
whose arithmetic status is inherited from the author's earlier work and is
\emph{not} proved here (\(\GAP\), declared); Proposition \ref{prop:cesaro}
guarantees that once supplied, it yields ``mean positivity of the symbol'',
delimiting the wall from inside: the wall is precisely the passage from mean
states to pure states.

\subsection{(c) What is genuinely new}
\begin{itemize}
\item The \emph{framework}: the Weil--Toeplitz algebra \(\W\), its corona
\(\Cor\), and the symbol \(\sigma(Q)\) (Definition \ref{def:field}, Lemmas
\ref{lem:symbol-norm}--\ref{lem:symbol-spectrum}), which realize the master
quantifier as \(c_0\) versus \(\ell^\infty\) (Corollary
\ref{cor:master-quantifier}).
\item The \emph{exact} characterization of robust defect finiteness (Theorem
\ref{thm:exact}, the (ii)\(\Rightarrow\)(i) direction), pinning the obstruction
to the single boundary point \(0\).
\item The \emph{trichotomy} (Theorems \ref{thm:fat}--\ref{thm:nogo}): off-critical
zeros stratify by visibility -- the thin-visible stratum on the symbol/compact
boundary, the fat stratum in the symbol, the sub-resolution stratum in the
ideal (Remark \ref{rem:trichotomy}) -- and \(m<\infty\) is a \(\Sigma_2\)
boundary event over the resolution filtration, certified by no fixed calendar.
\item That the \emph{realizability hypothesis} of \cite{P40} is replaced by the
strictly narrower, purely analytic gap of Hypothesis \ref{hyp:pw}
(the Paley--Wiener uniformity problem): test-accessibility is no longer the obstruction (Remark
\ref{rem:key}); what remains is a constant budget, open.
\item The exact logistic law (Theorem \ref{thm:logistic}) and the criterion
(Theorem \ref{thm:main}), whose RH branch is a theorem and whose \(\neg\)RH
branch is conditional on the Paley--Wiener uniformity problem.
\end{itemize}

\subsection{What remains open}
The computability of the symbol -- uniform pair-correlation asymptotics (Theorem
\ref{thm:symbolcompute}) -- is the old wall, intact. the Paley--Wiener uniformity problem (the
Paley--Wiener constant budget, Hypothesis \ref{hyp:pw} and Remark
\ref{rem:pw-rigidity}) is the single analytic gap of the \(\neg\)RH branch of
Theorem \ref{thm:main}. LP-\ref{conj:lp134}
(resolution repulsion) is a new, non-RH, falsifiable lemma; together with
replication (R) it would close the Fredholm half. The non-effectivity is
\(\Pi_1\), pure (Section \ref{sec:honest}(a)). We make the dependence on cited
inputs explicit.

\begin{remark}[cited hypotheses from earlier work on which the paper
depends]\label{rem:deps}
\textup{Theorem \ref{thm:main}}, the central equivalence, depends on:
\begin{enumerate}
\item the spectral side \eqref{eq:speca} with absolute convergence and the
factorization \(\hat h_{g,\varphi}=\half\hat g(\sigma)\hat\varphi(\tau)\), and
the free-orbit / hyperbolic-plane structure with multiplicity weights, from
\cite{P40} (and \cite{P35} for the realization \((\mathcal K,Q)\),
\(\kappa=2m\));
\item the autonomy of the value: \(Q_2(h,h')\) computable from primes without
zero positions, from \cite{P39,P35};
\item closure of the test class under multiplication by Mellin polynomials,
from \cite{P40};
\item under RH, diagonality of \eqref{eq:speca} (hence \(Q_2\geq0\)), from
\cite{P40};
\item the zero-density bound \(N(t+1)-N(t)\ll\log(2+t)\)
\cite[Thm.\ 5.8]{IwaniecKowalski} (published literature, not P-series);
\item for the \(\neg\)RH branch only: Hypothesis \ref{hyp:pw} (the Paley--Wiener uniformity problem),
\emph{open} -- not a cited input but a declared gap.
\end{enumerate}
\textup{The trichotomy (Section \ref{sec:trichotomy})} additionally depends on
the model axioms (S1)--(S3) of Definition \ref{def:model}, of which: the block
structure and the sub-resolution blindness (V1) are theorems of \cite{P40}; (S1)
(per-window pair-correlation boundedness) and (S3) (cross-window decay) are
declared \(\GAP\)s for the genuine Weil instance; (V2) carries the status of the
realizability lemma of \cite{P40}. \textup{Theorem \ref{thm:realiz}} uses items
(1)--(3),(5) and the open Hypothesis \ref{hyp:pw} (item (6), the Paley--Wiener uniformity problem),
but \emph{not} (S1)--(S3) and \emph{not} the realizability
lemma of \cite{P40} (Remark \ref{rem:key}). \textup{Proposition
\ref{prop:cesaro}} is unconditional given its stated density hypothesis, whose
arithmetic provenance (mean-value / large sieve) is a declared \(\GAP\).
\end{remark}

\appendix
\section{The constructive regime: Weil positivity for finite systems of primes}
\label{app:constructive}

The Weil--Toeplitz analysis of the body concerns the \emph{defect} --- the off-line
zeros and the negative directions they create. It is illuminating to record, by way
of contrast, the opposite extreme: the fully constructive regime, in which positivity
is immediate and the whole content lies in a single decidable discriminator. The
material of this appendix is the easy direction of Weil's criterion \cite{Weil1952},
recast categorically; we claim no novelty for the positivity itself and are explicit
about its standing.

\subsection{Pure Euler objects}
For a prime $p$ write $L_p=\log p$. The \emph{Euler object} of $p$ with parameter
$\alpha$ ($1\le|\alpha|\le p$) is the datum $F_{p,\alpha}$ whose local zeta factor is
$(1-\alpha p^{-s})^{-1}(1-(p/\bar\alpha)p^{-s})^{-1}$, with Weil functional
$W_{F_{p,\alpha}}$ assembled from the coefficients $c_m=\alpha^m+(p/\bar\alpha)^m$ of
its logarithmic derivative. The object is \emph{self-dual} (invariant under
$s\mapsto1-s$), and we say it satisfies the positivity property $H$ if its localized
Weil quadratic form $Q_{F_{p,\alpha}}(f,f)=W_{F_{p,\alpha}}(f\star\tilde f)\ge0$ for all
admissible $f$.

\begin{theorem}[Local purity is the exact discriminator]\label{thm:local-purity}
For $1\le|\alpha|\le p$,
\[
F_{p,\alpha}\text{ satisfies }H\quad\Longleftrightarrow\quad|\alpha|=\sqrt p
\quad\Longleftrightarrow\quad|c_m|\le2p^{m/2}\ \ \forall m\ge1 .
\]
\end{theorem}
\begin{proof}
If $|\alpha|=\sqrt p$ then $p/\bar\alpha=\alpha$, the two geometric progressions
coincide on $\Re s=\half$, and $c_m=2\alpha^m$ with $|c_m|=2p^{m/2}$. Summing the
single-factor Poisson identity over the doubled progression,
\[
Q_{F_{p,\alpha}}(f,f)=\sum_{\rho}\widehat f(\rho)\,\overline{\widehat f(\half+i\Im\rho)}
=2\sum_{\xi\in\Lambda_p}|\widehat F(\xi)|^2\ge0,
\]
since every zero $\rho$ has $\Re\rho=\half$ (so $1-\bar\rho=\rho$ and each term is a
squared modulus), where $\Lambda_p=\{(\arg\alpha+2\pi k)/L_p:k\in\mathbb Z\}$.
Conversely, if $|\alpha|=p^{1/2+\delta'}$ with $\delta'\ne0$, the reverse triangle
inequality gives $|c_m|\ge p^{m/2}(p^{m\delta'}-p^{-m\delta'})>2p^{m/2}$ for large $m$,
the zeros leave the line, and $H$ fails. Equality $|\alpha|=\sqrt p$ is exactly
$|c_m|=2p^{m/2}$.
\end{proof}

\subsection{Additivity: the cross terms vanish exactly}
\begin{proposition}[Convolution does not mix Weil data]\label{prop:crosszero}
For Euler objects $X_1,X_2$ of distinct primes, the Weil functional of the product is
additive, $W_{X_1\star X_2}=W_{X_1}+W_{X_2}$, with cross terms \emph{exactly zero}.
Consequently a product of pure Euler objects over any finite set of primes satisfies
$H$ if and only if each factor does, i.e.\ if and only if each is pure
\textnormal{(}$|\alpha_p|=\sqrt p$\textnormal{)}.
\end{proposition}
\begin{proof}
The logarithmic derivative of a product is the sum of the logarithmic derivatives, and
the Weil functional is linear in it; the cross terms are integrals of
$e^{i(\log p)\xi}\,\overline{e^{i(\log q)\xi}}$ against the test data, which vanish
because $\log p\ne\log q$ places the two progressions on disjoint frequency lattices.
Hence positivity is per-prime, and Theorem~\ref{thm:local-purity} applies factorwise.
\end{proof}

\subsection{Honest scope}
Proposition~\ref{prop:crosszero} with Theorem~\ref{thm:local-purity} is exactly the
\emph{easy} direction of Weil's criterion: purity places every zero on $\Re s=\half$
\emph{by hypothesis}, and positivity follows term by term, with the additivity of
logarithmic derivatives (classical since Euler) making the finite system no harder than
one prime. We claim no novelty for the positivity. Two structural points are worth
recording. First, purity is the \emph{exact, decidable} discriminator: $H$ holds if and
only if $|\alpha_p|=\sqrt p$ for every $p$ in the system, a sharp local condition.
Second, the cross-prime interference, which one might fear, is located precisely: it is
absent from the \emph{sign} (the cross terms are identically zero) and present only in
the spectral structure that a genuine global object would carry. The finite Euler
product is \emph{not} $\zeta$: the constructive positivity survives every finite
truncation but the spectrality does not, and $\zeta$ is not the limit of its Euler
factors in the relevant sense. The contrast with the body is the point: in the
constructive regime purity is assumed and positivity is free, whereas the difficulty
addressed in the main text is precisely the absence of an a priori purity.

\section*{Correction history}
\emph{June 2026.} Two adversarial audits mandated the following corrections, all applied in this version. A
false victory being worse than a failure, we record them explicitly.
\begin{itemize}
\item \textbf{(Critical.)} The \(\neg\)RH branch of the
logistic criterion was previously claimed as a theorem over cited inputs. Its
proof contained a constant treated as constant while depending on the parameter
sent to infinity: the tail budget of Theorem \ref{thm:realiz}, Step 3, used
\(M\)-dependent Paley--Wiener constants whose uniformity is not only unproved
but \emph{false} for every fixed bump (Paley--Wiener rigidity, Remark
\ref{rem:pw-rigidity}). Status downgraded: the branch is now \emph{conditional}
on Hypothesis \ref{hyp:pw} (the Paley--Wiener uniformity problem). Abstract, introduction, Theorem
\ref{thm:main} and Section \ref{sec:honest} reworded accordingly. The statement
itself is not refuted; the RH branch (Proposition \ref{prop:RHbranch}) and the
exact logistic algebra (Theorem \ref{thm:logistic}) survive the audit intact.
\item \textbf{(Correction.)} The same proof failed to annihilate the
\emph{resonant pairs} (both members in \(S\), \(\sigma\in\{\sigma_0,\sigma_0'\}\),
off the target orbit), which exist under general \(\neg\)RH configurations and
contribute terms of uncontrolled sign. Repaired here (repair validated by the
audit) by Vandermonde annihilation on the \(\tau\) side over \(Z_\tau\),
enlarging \(D\); the enlargement feeds back into the constant budget, so the
two gaps must be closed jointly within the Paley--Wiener uniformity problem.
\item \textbf{(Correction.)} The closed logistic formula in
\(\sigma_{\mathrm{loc}}\) (the former version of Proposition
\ref{prop:loc-logistic}) was false when the
winning Gram is degenerate (two incompatible normalizations of \(x_0(N)\):
the radical bug). Repaired via
\(\sigma_{\mathrm{base}}^{\,\mathcal U}\in[\sigma_{\mathrm{loc}},
2\sigma_{\mathrm{loc}}]\cap[0,\half]\) (Remark \ref{rem:radical-bug}); the
\(0\)-versus-\(\half\) dichotomy and the \(\mathcal U\)-independence of the
limit survive; intermediate floors are no longer claimed canonical.
\item \textbf{(adversarial audit)} ``\(\negind\) computed
exactly'' weakened to semidecidability by interval arithmetic (Proposition
\ref{prop:loc-welldef}); mesh constant corrected to \(b\leq\mu/(4\sqrt2\,L)\)
and ``computable \(N_0\)'' to ``explicit \(N_0\)'' (Theorem \ref{thm:mesh});
the unconditional ``\(\to0\) as \(M\to\infty\)'' claim removed from Lemma
\ref{lem:tail}; the attribution of the global-signature collapse corrected
(Remark \ref{rem:collapse}).
\item \textbf{(adversarial audit)} The noisy-case exception clause of Theorem
\ref{thm:strict} (``up to windows with \(\vartheta_n\leq\omega_n\)
eventually'') was \emph{false} (counterexample: \(\omega_n=1/n\),
\(\vartheta_n=1/(2n)\)). The correction \emph{strengthens} the theorem: strict
symbol positivity \(\iff\) finitely many visible windows, with no exception
(Remark \ref{rem:strict-corrected}). Downstream results unaffected.
\item \textbf{(adversarial audit)} ``Continuous field'' replaced by
``bounded field over the discrete base \(\N\)'' (over a discrete base the
Dixmier continuity condition is vacuous; no proof used it); the global slogan
of Remark \ref{rem:trichotomy} stratified (only the thin-visible stratum lives
on the symbol/compact boundary); a typographic inequality chain in Lemma
\ref{lem:window-spec} fixed; Theorem \ref{thm:thin}(i) qualified by
\(\vartheta_n>\omega_n\) i.o.\ for the negativity claim. Theorem
\ref{thm:exact} and Lemmas \ref{lem:symbol-norm}--\ref{lem:symbol-spectrum}
passed the audit without correction.
\end{itemize}

\section*{Acknowledgements}
This work builds on the author's earlier works \cite{P35,P39,P40}, cited as prior
work.

\begin{thebibliography}{99}

\bibitem{Bhatia}
R.~Bhatia,
\emph{Matrix Analysis},
Graduate Texts in Mathematics \textbf{169}, Springer, 1997.

\bibitem{BottcherSilbermann}
A.~B\"ottcher and B.~Silbermann,
\emph{Analysis of Toeplitz Operators},
2nd ed., Springer, 2006.

\bibitem{Gantmacher}
F.~R.~Gantmacher,
\emph{The Theory of Matrices}, Vol.~2,
Chelsea, 1959.

\bibitem{IwaniecKowalski}
H.~Iwaniec and E.~Kowalski,
\emph{Analytic Number Theory},
American Mathematical Society Colloquium Publications \textbf{53}, AMS, 2004.

\bibitem{Kato}
T.~Kato,
\emph{Perturbation Theory for Linear Operators},
2nd ed., Springer, 1976.

\bibitem{RabinovichRochSilbermann}
V.~Rabinovich, S.~Roch, and B.~Silbermann,
\emph{Limit Operators and Their Applications in Operator Theory},
Operator Theory: Advances and Applications \textbf{150}, Birkh\"auser, 2004.

\bibitem{ReedSimon}
M.~Reed and B.~Simon,
\emph{Methods of Modern Mathematical Physics, IV: Analysis of Operators},
Academic Press, 1978.

\bibitem{Robinson}
A.~Robinson,
\emph{Non-standard Analysis},
North-Holland, 1966.

\bibitem{Weil1952}
A.~Weil,
\emph{Sur les ``formules explicites'' de la th\'eorie des nombres premiers},
Comm.\ S\'em.\ Math.\ Lund (1952), 252--265.

\bibitem{P35} D.~A.~Trejo Pizzo, \emph{Pontryagin spectral index of the Weil form and the Connes scaling operator: a bridge theorem for the Riemann Hypothesis}, preprint, 2026.

\bibitem{P39} D.~A.~Trejo Pizzo, \emph{A trace and measure criterion for the Riemann Hypothesis from the zeta spectral triple}, preprint, 2026.

\bibitem{P40} D.~A.~Trejo Pizzo, \emph{The localized Weil form and its Kre\u\i n--Pontryagin index: localization, obstruction, and rigidity}, preprint, 2026.

\end{thebibliography}

\end{document}
